\subsection{Gamma-dominant and root-coalescence faces}
\label{subsec:deg-gamma-dominant}

In the chart \eqref{deg:gamma-chart}, negative \(\gamma\) leaves two
separated incident roots as \(k\to0\). For positive \(\gamma\), the
positive roots may coalesce and the negative root escapes. These two
geometries require different comparisons: a physical endpoint power law
on the negative face, and weighted curvature on the positive face.
We first treat the negative face and the positive far-root face, then
the bounded positive face. In each case the displacement is the original
\(F(x)=x-m\Pi_{\theta,m}(x)\) of \eqref{eq:return}, with
\(x=h-a_\theta\). Fix the original passive compact set, nested windows
and their reserves before the regional cuts. Every carrier below is
\[
 I_\theta=\{x:a_\theta+x\in J_\theta^{\rm out},\quad
                  0<H(x):=z_\theta(\gamma;a_\theta+x)<H_{\rm cap}\},
\]
where \(0<H_{\rm cap}<z_*\) is a positive unscaled cap chosen before
the coefficient and multiplier bounds. Thus every set retains the original
input window, inverse output window, stop and speed clips, and all
intervening nonroot inputs. Phase derivatives fix every coefficient and
\(m\); a stopped finite channel contributes no cycle.

\begin{proposition}[Negative quadratic face]
\label{deg:prop-negative-gamma}
Fix \(\Pi_{\max}<\infty\) and a sufficiently small routing bound
\(k_\flat\) from Lemma~\ref{deg:lem-small-source-routing}.
For \(\gamma=-g<0\), \(0<k\le2k_\flat\), and
\(0\le\pi\le\Pi_{\max}\), independently small \(g\) and \(m\)
give the physical endpoint law \eqref{deg:negative-gamma-physical-power}.
The endpoint, core and high intervals cover the original \(I_\theta\),
with strict estimate overlaps and positive dividing heights retained.
Every critical point of the same \(F\) is simple and has \(F''>0\).
Consequently the whole coefficient fibre has at most one critical point
and at most two isolated displacement zeros, counted with multiplicity.
\end{proposition}

\begin{proof}
We first match the two root passages on fixed normal sections, retaining
the middle multiplier in their logarithmic actions. The exterior maps
then turn this match into the physical power law. A comparison orbit
joins its endpoint interval to the bounded core and high-height interval,
so the final critical-point argument applies to one connected carrier.

Put \(\epsilon=gk\). The exact dictionary is
\(v-\gamma=\epsilon\xi\), \(z=\epsilon^2\eta=s^2Y_0\),
\(\eta=gY_0\), \(s=g^{3/2}k\). After positive time scaling,
\[
 \dot\xi=-gP-Q\eta,\qquad \dot\eta=-\eta(H_0+H_1\eta),
\]
\[
 \begin{gathered}
 P=1-\pi\xi-3\xi^2+k\xi^3,\quad
 Q=1+Ag(1-k\xi)+Cg^2(1-k\xi)^2,\\
 H_0=(1-2g)\xi+kg\xi^2,\quad H_1=-Ckg^2(1-k\xi).
 \end{gathered}
\]
Its incident roots \(a>0,-b<0\), their distances from zero,
\(|P'(a)|,P'(-b)\), and normal rates are positive units on the
closed \((\pi,k)\) block. The third root lies beyond a fixed large
cut \(L\). Choose independent fixed normal sections and inner weak
cuts at the two incident roots. The intrinsic construction
\eqref{eq:saddle-normal-time}--\eqref{eq:exact-saddle-invariant}
has kernels \(C(1+T)^Ne^{-(c-Cg)T}\); four differentiated integral
equations give uniformly bounded strong projections and normal coordinates.
With inward footpoint distances \(t,r>0\), put
\(A_+=-H_0(a)/P'(a)\), \(A_-=-H_0(-b)/P'(-b)\),
\(\nu=A_+/A_-\). These are positive units. The anchored normal-time
actions satisfy
\begin{equation}\label{deg:negative-gamma-actions}
 \begin{gathered}
 J_+(t)=-A_+\log t+j_+(t),\quad J_-(r)=-A_-\log r+j_-(r),\\
 |(\delta\partial_\delta)^j(j_\pm(\delta)-j_\pm(0))|\le C\delta,
 \quad j\le3.
 \end{gathered}
\end{equation}
Indeed, in the intrinsic equations
\(\dot\Xi=-\rho\Xi,\ \dot Z=Z,\ \rho=g/A_\pm\),
\(|\Xi|Z^\rho\) is constant and
\(J=-(g/\rho)\log(|\Xi_n|/|\Xi_m(0)|)\); its axis derivative
is \(H_0/P\). Pole subtraction gives the displayed smooth remainders.
On each actual cut write
\(|\Xi_n|=\delta a_n(\delta)\), \(Z_n=b_n(\delta)\),
\(Z_m=z_0wU_m(w)\). The positive units have uniform \(C^3\) bounds;
Taylor division uses the intrinsic fourth derivative. The exact passage is
\[
 w=w_0e^{\Psi_g(w)},\qquad
 w_0=(b_n/z_0)e^{-J/g},\qquad
 \Psi_g=\rho^{-1}\log(|\Xi_m(0)|/|\Xi_m(w)|)-\log U_m(w).
\]
Here \(|\Psi_g|\le Cw/g\), \(|\partial_w^j\Psi_g|\le C/g\),
\(j\le3\). Choose each fixed collar before \(g\), so that
\(w_0\le\delta^{A_{\min}/(2g)}\). Then
\(1-w\Psi_g'(w)\ge1/2\) and \(w\asymp w_0\). Implicit Euler
differentiation gives
\(|(\delta\partial_\delta)^j\log(w/w_0)|\le Cg^{-(j+1)}w\),
\(j\le3\), hence ordinary \(C^3\) extension to zero. Thus,
independently on both sections, \(w_\pm=u_\pm e^{-J_\pm/g}\),
where \(u_\pm\) and their inverses are uniform \(C^3\) units.

Take \(w\) to be the raw height \(\eta\) on each fixed inner
weak cut, where the normal restriction is \(Z_m=z_0wU_m(w)\).
On the core \(P\ge p_c>0\),
the scalar equation is
\(f=\eta(H_0+H_1\eta)/(gP+Q\eta)\).
Choose an actual input chord \(0\le w\le e^{-K_0/g}\), with
\(K_0\) fixed large enough that a first-exit comparison keeps its
whole passage and output inside the two weak-cut buffers.
Here \(|f_{\eta^h}|\le C_hg^{-h}\). The initial-height rows
\(V_j=\partial_w^j\eta\), through order five, satisfy
\[
 V_j'=f_\eta V_j+
 \sum_{h=2}^j f_{\eta^h}\mathcal B_{j,h}(V_1,\ldots,V_{j-h+1}),
\]
where \(\mathcal B_{j,h}\) are the partial Bell polynomials.
The linear kernel is only \(O(g^{-1})\); successive variation of
constants puts the higher inverse powers in the forcing and gives
\mbox{\(\|M_g\|_{C^5}+\mu_g^{-1}\le e^{K/g}\)}, after enlarging \(K\).
For the two fixed physical cuts the linear multiplier is exactly
\begin{equation}\label{deg:negative-core-multiplier}
 g\log\mu_g=\int_{\xi_{i,+}}^{\xi_{i,-}}\frac{H_0}{P}\,d\xi .
\end{equation}
Taylor's integral division
\(N_g(w)=\mu_g^{-1}\int_0^1(1-t)M_g''(tw)\,dt\)
therefore gives the actual middle map
\begin{equation}\label{deg:negative-gamma-middle}
 M_g(w)=\mu_gw(1+wN_g(w)),\qquad
 \mu_g^{\pm1}+\|N_g\|_{C^3}\le e^{K/g}.
\end{equation}
Choose \(t_0\) before \(g\), with its action exceeding all these fixed
exponents. The weak input then lies on the actual positive core chord,
and \(g(t\partial_t)^j\log(1+wN_g(w))=O(t^N)\), \(j\le3\),
for any fixed \(N\). Matching the actual heights, while retaining
\(g\log\mu_g\), gives
\begin{equation}\label{deg:negative-gamma-matching}
 \begin{gathered}
 A_-\log r-A_+\log t=C_\theta+e_+(t)+e_-(r),\qquad |C_\theta|\le C,\\
 |(t\partial_t)^je_+|\le C(t+t^N),\quad
 |(r\partial_r)^je_-|\le Cr,\quad j\le3.
 \end{gathered}
\end{equation}
The negative height is strictly monotone from zero, and
\(A_--r\partial_re_->0\). Reducing \(t_0\) admits a unique actual
output in its independent fixed negative collar. Consequently
\begin{equation}\label{deg:negative-gamma-power}
 \begin{gathered}
 r=A_\theta t^\nu e^{E_\theta(t)},\qquad c\le A_\theta,A_\theta^{-1}\le C,\\
 |(t\partial_t)^jE_\theta|\le Ct^\beta,\quad j\le3,
 \qquad0<\beta<\min(1,\inf\nu).
 \end{gathered}
\end{equation}

Continue both sections to \(\pm L\). Their fixed-normal fast graphs
have limit \(\eta=(\xi^2-u^2)/2\); at the cuts their heights,
denominators and gap derivatives are units. At a fixed small raw cut
\(d\), differences of the bounded \(L\)-labels are
\(O(\epsilon^{2-\lambda_\gamma})\), while the reference graph has
\(z/d^2=1/2+O(d)+o(1)\). The raw corridor first variation is
\(C(d_1/d)^{\lambda_\gamma}\) times a unit. Choosing \(d\) first
and then \(g\) therefore propagates the reference to the original
organizing-arc windows uniformly on the retained passive compact set.
This admits the whole outward chord, including its limiting zero end,
before applying its ordinary observers; it imposes no condition on
\(g\log k\). The exact tail calculation
\eqref{ann:tail-first-derivative}--\eqref{eq:N-bound} yields
\[
 d_\pm=\epsilon^{2-\lambda_\gamma}U_\pm,\qquad
 \lambda_\gamma=-12g/(1-2g),\qquad
 c\le U_\pm,U_\pm^{-1}\le C.
\]
Its first remainder is integrable with majorant
\(C(\epsilon+g\pi R^{-2}+R^{-3+\zeta})\), \(0<\zeta<1\);
the higher normalized kernels give ordinary \(C^3\) embeddings.
Only \(d_-\) includes the divided closing, whose first derivative is
\(f'(md)d'\). With exact centers and slopes retained separately,
\[
 x=x_{\rm st}+d_+\phi_+(t),\qquad
 \Pi(x)=y_{\rm st}+d_-\phi_-(r(t)),\qquad
 \phi_\pm(0)=0,\quad\phi_\pm'(0)=1.
\]
These maps and their inverses are uniformly \(C^3\). Differentiating
this identity, and putting \(q=(x-x_{\rm st})/d_+=\phi_+(t)\), gives
\begin{equation}\label{deg:negative-gamma-physical-power}
 \begin{gathered}
 \Pi'(x)=B_{\theta,m}q^{\nu-1}e^{R(q)},\qquad B_{\theta,m}=(d_-/d_+)A_\theta\nu,\\
 c\le B_{\theta,m},B_{\theta,m}^{-1}\le C,\qquad
 |(q\partial_q)^jR|\le Cq^\beta,\quad j\le2.
 \end{gathered}
\end{equation}
In fact
\(R(\phi_+(t))=E+\log(1+tE'/\nu)+\log\phi_-'(r)
-\log\phi_+'(t)+(\nu-1)\log(t/\phi_+(t))\).
The third Euler row in \eqref{deg:negative-gamma-power} controls the two
Euler derivatives of this expression. The absolute amplitude
\(B_{\theta,m}d_+^{1-\nu}\) need not be a unit.

Choose the comparison \(t_c=t_0/2\); its matched \(r_c\) lies in a
fixed positive interval. On the two outer pieces beyond
\(a-t_c/2\) and \(-b+r_{\min}/2\), its fast height has a fixed
positive margin. The intervening piece has \(P\ge p_c\).
Choose a fixed upper central shot \(\eta_{\rm core}>0\); its limiting
graph is \(\eta_{\rm core}+\xi^2/2\). Both bounding shots have
buffered outward admission as above. Scalar order and \(Q>0\) admit
every intervening shot \(\eta_c\le\eta_0\le\eta_{\rm core}\),
including nonroots, and their raw central speed is at most
\(\epsilon^2(g+Q(0)\eta_{\rm core})<d_{\rm out}\).
For \(D=gP+Q\eta\) and \(f=\eta(H_0+H_1\eta)/D\),
use \eqref{deg:scalar-raw-cocycle}, with residual
\([gP'+(Q'+H_1)\eta]/D\).
On the middle piece the last fraction is bounded by
\(C+Cgk\); on the outer pieces the comparison height bounds
\(D\) below. The endpoint quotients are units, so integration bounds
the core first variation and its inverse. The two outward derivatives
are \(\epsilon^{2-\lambda_\gamma}\) times units on this entire
chord. Their ratio, regular bridges and closing give \(0<\Pi'\le C_c\).
The limiting \(t=0\) orbit stops, whereas every \(t>0\) passes;
Lemma~\ref{lem:loc-through} identifies \(x_{\rm st}\) with the
excluded mother endpoint. The endpoint extends through \(t_0\),
strictly beyond the comparison.

The comparison action gives \(\eta_c\le Cg^{-N}e^{-c/g}\). The
signed high threshold and the ordered final shrink are
\begin{equation}\label{deg:negative-gamma-cover}
 H_H=\epsilon^2gK_H(1+g\pi^2),\qquad
 \eta_c<\eta_H:=gK_H(1+g\pi^2)<\eta_{\rm core}/2,\qquad
 \epsilon^2\eta_{\rm core}<H_{\rm cap}/2.
\end{equation}
Hence the endpoint, closed core and signed high packages strictly overlap
and cover \(I_\theta\) on the original height axis. The equality
\(H=H_H\) belongs to the core; original clips are imposed on all
three. Choose \(m\) last so \(m\max(C_c,C_H)<1/2\).
It remains to control critical points in the endpoint interval,
including parameters for which the exponent ratio approaches one.
On the endpoint,
\(F'=1-e^L\), \(L=\log(mB)+(\nu-1)\log q+R(q)\).
If \(\nu\ge1\), then \(L<0\). Otherwise, for \(z=1-\nu>0\),
a critical point satisfies \(z\log(1/q)\ge M:=\log(1/m)-C\), so
\[
 qL_q=-z+O(q^\beta)\le-z+Ce^{-\beta M/z}<-z/2.
\]
This choice is uniform as \(z\downarrow0\). Since \(q_x=1/d_+>0\),
every actual critical point has \(F''>0\). The core and high add none;
Lemma~\ref{sharp:upper-critical-count} gives the whole-carrier count.
\end{proof}

\begin{proposition}[The positive far-root face]
\label{deg:prop-positive-far}
For \(\gamma=g>0\), there are independent bounds
\begin{equation}\label{deg:positive-far-range}
 0<\lambda=\pi^{-2}\le\lambda_0,\qquad
 0<q=k\pi\le q_0,\qquad 0<\tau=g\pi^2=p^2\le\tau_0
\end{equation}
such that the actual lower interval has strictly negative completed-map
curvature. Its bounded core and signed high complement strictly overlap
and have \(F'>1/2\) after one multiplier shrink. They cover the
original raw-capped \(I_\theta\), including all nonroot inputs and
original clips. The same physical \(F\) has at most one critical point
and at most two isolated zeros, counted with multiplicity, on every fibre.
\end{proposition}

\begin{proof}
The negative incident root now escapes in the scaled coordinate.
We estimate its outward map first below one third of the root radius and then
in a root collar rescaled by that radius. The exact corridor exponent
keeps both estimates uniform as the root escapes; their action weights
make the resulting section correction small in physical curvature.

Put \(\chi=q\lambda^2\), \(a=s/p=\tau\chi\),
\(g=\tau\lambda\), \(s=\tau^{3/2}\chi\),
\(v=g+aU\), \(z=a^2Z\); at the center \(Z=\tau Y_0\).
After positive time scaling the exact field is
\[
 U'=-\tau P_+(U)-Q_+(U)Z,\qquad
 Z'=-Z\{UB_+(U)+Ca(g+aU)Z\},
\]
\begin{equation}\label{deg:far-polynomials}
 \begin{gathered}
 P_+=1-U+3\lambda U^2+\chi U^3,\qquad
 P_-=1+T+3\lambda T^2-\chi T^3,\\
 B_+(U)=1+2g+aU,\quad B_-(T)=1+2g-aT,\\
 Q_+(U)=\ell(g+aU),\quad Q_-(T)=\ell(g-aT),\quad
 \ell(v)=1-Av+Cv^2.
 \end{gathered}
\end{equation}
The moving lips are \(Z=\tau\ell_0P_\pm\), with labels
\(b_+=-U^2/2,b_-=-T^2/2\). The source quartet is
\((-\tau,\tau,\tau\lambda,\tau\chi)\); reflection at a fixed
negative base \(T_b\) gives
\((-\tau/T_b^2,-\tau/T_b,\tau\lambda,-\tau\chi T_b)\)
without changing the positive-time orbit order.
The first positive root is \(A_+=1+O(\lambda+\chi)\).
Writing \(w_*=q\lambda T_*\), the unique negative-side root satisfies
\begin{equation}\label{deg:far-roots}
 w_*=3+q/w_*+q^2\lambda/w_*^2,\quad 3<w_*<4,\qquad
 T_*\asymp(q\lambda)^{-1},\quad \chi T_*\le4\lambda.
\end{equation}
Its slope is \(c_-=2+3\lambda T_*+3/T_*\). The action residues
\(r_+=A_+B_+(A_+)/[-P_+'(A_+)]\asymp1\) and
\(r_-=T_*B_-(T_*)/c_-\asymp\lambda^{-1}\) are retained exactly.
Define
\[
 \begin{gathered}
 J_+(U)=\int_0^U\frac{xB_+}{P_+}\,dx,\qquad
 A_-(T)=\int_0^T\frac{xB_-}{P_-}\,dx,\\
 \mathcal D_\pm=\frac{P_\pm}{xB_\pm}\partial_x,\quad
 H_\pm=\frac{P_\pm'}{xB_\pm},\quad h_\pm=\frac{P_\pm}{x^2B_\pm}.
 \end{gathered}
\]
With \(R=T^{-1},d=\chi T,E=T\partial_T\),
\(H_-=(R+6\lambda-3d)/B_-\),
\(h_-=(R^2+R+3\lambda-d)/B_-\), \(\mathcal D_-=h_-E\).
Repeated differentiation gives
\begin{equation}\label{deg:far-action-bound}
 \frac{|\mathcal D_-^jH_-|}{h_-}
 \le C_j(R+\lambda)^{j-1}(R+d),\qquad
 \int_{T_b}^{T_*}|\mathcal D_-^jH_-|\,dA_-
 \le C_j(T_b^{-1}+\chi T_*),\quad j=1,2,3.
\end{equation}
Here \(dA_-=d\log T/h_-\); the retained factor \(R+d\) avoids
any \(\lambda\log T_*\) loss. Also \(\tau H_-\le1/4\).
The reflected coefficient \(\varpi_-=-C\tau a(g-aT)P_-/(TB_-)\)
and its first four radial rows are \(O(\tau^3\lambda^3)\), while
those of \(Q_--1\) are \(O(g)\). Choose \(\ell_0\) with
\(\ell_0\max(|Q_-|,|\varpi_-|)\le1/4\). The positive root
collar has the same oriented margins on its fixed bounded cuts.

Apply Lemma~\ref{deg:lem-stable-ratio-ledger} to
\(\vartheta=Z/(\tau P)\), \(q_a=(A(x_l)-A(x))/\tau\), with
\(D=\mathcal D_l\) on either entrance, \(\kappa=\tau\), \(\omega=\varpi\),
\(\rho=\tau\), and \(\alpha=0\). At fixed clock, endpoint
transport is \(\mathcal D\) at the running point; the coefficient
rows follow from \eqref{deg:far-action-bound} and the bounded positive
collar. The margins above give the two negative diagonals, and
\(\partial_{q_a}=-\tau\mathcal D\) supplies the clock rows.
Thus \(\vartheta\le\ell_0e^{-q_a/4}\), and the exact memory is
\[
 \vartheta=\ell_0\frac{P(x_l)}{P(x)}e^{-q_a+V},\qquad
 V=\int_0^{q_a}\frac{(Q-\varpi)\vartheta}{1+Q\vartheta}\,dq,
\]
with three bounded action rows. In particular the initial polynomial
factor is retained through the long negative flight.
Fix inner cuts \(U_i>0\), \(0<T_i<T_b\) on a positive core,
with \(A_-(T_b)-A_-(T_i)\ge c_i>0\), and put
\(J_i=J_+(U_i)\), \(A_i=A_-(T_i)\).
At either fixed inner cut the terminal clock is
\(q_i=(A(x_l)-A(x_i))/\tau\), with
\(\mathcal D_lq_i=1/\tau\) and
\(\mathcal D_l^2q_i=\mathcal D_l^3q_i=0\).
All three total rows in \eqref{deg:stable-ratio-terminal-rows}
are bounded: their terminal densities decay as \(e^{-c/\tau}\),
absorbing the factors through \(\tau^{-3}\).
Writing \(V_\pm\) for the complete terminal integrals, the actual
raw weak heights are
\[
 \begin{aligned}
 w_+(U)&=\tau\ell_0P_+(U)e^{-[J_+(U)-J_i]/\tau+V_+(U)},\\
 w_-(T)&=\tau\ell_0P_-(T)e^{-[A_-(T)-A_i]/\tau+V_-(T)}.
 \end{aligned}
\]
The same five-row scalar core construction, with \(g\) replaced by
\(\tau\), gives
\(M(w)=\mu_\theta w(1+wN_\theta(w))\) on its actual weak chord,
with \(\mu_\theta^{\pm1}+\|N_\theta\|_{C^3}\le e^{K/\tau}\).
Its exact first multiplier has
\(\log\mu_\theta=(A_i-J_i)/\tau\).
Fix the cuts, large radial cut and normal height first; enlarge \(K\)
to include the fixed action-bracket reserves, then choose
\(X_0>0\) with \(J_+(A_+-X_0)-J_i>4K+1\), and
\(0<X_1<X_0\). The actual weak memory and three action rows are
\(O(\tau^{-N}e^{-(3K+1)/\tau})\) on \(A_+-X_0\le U<A_+\).
The exact equality \(M(w_+(U))=w_-(T)\) now gives
\begin{equation}\label{deg:far-action}
 J_+(U)-A_-(T)=\tau\{\log P_+(U)-\log P_-(T)+R_0(U,T)\},
 \quad |R_0|+\sum_{1\le i+j\le3}|\mathcal D_+^i\mathcal D_-^jR_0|\le C.
\end{equation}
Here the remainder has the separated extension
\begin{equation}\label{deg:far-separated-remainder}
 R_0(U,T)=V_+(U)-V_-(T)
       +\log\{1+w_+(U)N_\theta(w_+(U))\}.
\end{equation}
It is defined on the product of the actual positive weak chord and
negative inward domain. Its mixed rows with \(i,j>0\) are exactly zero;
the pure rows through order three are those just estimated.
Only then is it restricted to the actual center-image intersection.
Scalar order gives its unique matched
\(T_b\le T<T_*\); at both comparison labels \(A_+-X_j\),
fixed action brackets and \(A_-(T)\to T-\log(1+T)\) put \(T_j\)
in fixed positive compact intervals. Differentiation yields
\[
 T_U=\frac{UB_+P_-}{TB_-P_+}(1+\varrho),\qquad
 1+\varrho=\frac{1-\tau H_+-\tau\mathcal D_+R_0}
                   {1-\tau H_-+\tau\mathcal D_-R_0},\qquad
 |\mathcal D_+^i\mathcal D_-^j\varrho|\le C\tau\quad(i+j\le2).
\]

Let \(\Gamma_\pm(b_\pm)=d_\pm+c_\pm\phi_\pm(b_\pm)\) be the
actual raw embeddings, including the regular bridges and output closing;
the positive center is measured in \(x=h-a_\theta\).
With \(S(b_+)=-T^2/2\),
\(\widehat\Pi=\phi_-\circ S\circ\phi_+^{-1}\) and
\(\Pi(d_++c_+w)=d_-+c_-\widehat\Pi(w)\).
Put \(L_\pm=\mathcal D_\pm\log\Gamma_\pm'\), where the primes
refer to their own \(b_\pm\). Apply \eqref{deg:weighted-exact-curvature-chain} with
\(u=U,v=T,W=(U+T)/(UT)\), retaining both section terms.
In this weight the polynomial part is
\(3\chi UT-1\), apart from \(O(\tau\lambda)\).
The relative correction is \(O(\tau)\); the section term has bounded
weight \(UB_+T/(U+T)\). We now estimate its actual log-slopes.

At the positive root, pole subtraction and the intrinsic normal clock
retain the embedding power \((A_+-U)^{1+\alpha_+}\),
\(\alpha_+=\tau/r_+\). The normal distance kernel and its first
three weighted rows have majorant \(C(W+\tau u)\), whose normal-clock
integral is \(O(\sigma+\tau(1+|\log\tau|))\).
Its fixed normal-to-raw ordinary \(C^3\) observer therefore gives
\[
 |L_+|\le C_0\sigma+C_\sigma X_0+
                         C_\sigma\tau(1+|\log\tau|)^3.
\]
Choose \(\sigma\), then \(X_0\), to make the fixed terms small;
reducing \(X_0\) strengthens the earlier action gap.

For the negative outward graph put
\(D=\tau P_-+Q_-Z\), \(n=RB_-\),
\(\zeta=-Ca(g-aR)\), \(f=Z(n+\zeta Z)/D\).
It starts at the true moving value \(Z_l(T)=\tau\ell_0P_-(T)\)
and ends at the fixed raw port \(R_{\rm out}=\delta_-/a\).
The original observer \(\psi\) is composed only after raw admission.
Use \eqref{deg:scalar-polynomial-lip-slope} with
\(M=\tau P_-\), \(x_l=T\), \(\sigma=1\), and
\(\chi(Z)=\psi(a^2Z)\). It gives
\[
 \Gamma_-'(b_-)=a^2\psi'(a^2Z_T)\frac{Z_T}{D_T}B_-(T)
 [1+\ell_0\varpi_--\tau H_-(1+\ell_0Q_-)]e^{I(T)},
\]
where \(I=\int_T^{R_{\rm out}}[\tau P_-'+(Q_-'+\zeta)Z]/D\,dR\).
This retains the exact moving initial factor and first gain.
The two logarithmic-slope rows use only the ordinary third observer
derivative, by \eqref{phys:eq-observer-log}.

For \(T_b\le T\le T_*/3\), apply
Lemma~\ref{deg:lem-spectral-exterior} in the raw coordinate \(x=aR\),
with starting abscissa \(aT\), terminal cut \(\delta_-\), and
\[
 t=-1,\quad (\mu_2,\mu_1,\mu_0)=(3g,\tau a,\tau a^2),\quad
 b=2g,\quad v=g,\quad \sigma=a.
\]
Its input label is exactly \(-T^2/2\). With the comparison mass
\(\kappa=\tau/T^2+\tau/T+g+aT\), the root bound \(aT_*<4g\)
gives \(m(x)\asymp\kappa\) on \([aT,2aT]\), and all coefficient
bounds in \eqref{deg:spectral-exterior-hypotheses} hold.
The same margin holds for \(T<7T_*/20\), providing a strict overlap
with the root chart beginning at \(T_*/3\).
Here \(\Theta\le\kappa\le C\tau\), the terminal coefficient is
bounded, and \((aT/\delta_-)^{2-\nu}\le C\tau\), where
\(\nu=12g/(1+2g)\) is the lemma's exact spectral exponent.
Thus, for \(E=T\partial_T\), its two rows give
\[
 |E\log\Gamma_-'|+|E^2\log\Gamma_-'|
                       \le C\tau(1+|\log\tau|)^3.
\]
The observer is used on its original admitted image; neither \(a\)
nor the comparison mass is differentiated as a coefficient.

For the rescaled root range below, the corridor from \(R_c=LT_*\) retains
\(y=Z/R^2\), \(Q_g=\ell(g)\),
\(y_g=(1-4g)/(2Q_g)\), \(\nu=12g/(1+2g)\) exactly.
For \(m_R=\tau/R^2+\tau/R+3g-aR\),
\(B=1+2g-aR\), \(Q=\ell(g-aR)\),
\(\zeta_R=-CaR(g-aR)\), \(D_y=m_R+Qy\), the exact equations are
\[
 R y_R=y\left(\frac{B+\zeta_Ry}{D_y}-2\right),\qquad
 R(y-y_g)_R=-\frac{2Q_gy}{D_y}(y-y_g)+S,
\]
\[
 S=\frac y{D_y}\{(B-1-2g)-2(m_R-3g)+
                                      [\zeta_R-2(Q-Q_g)]y\}.
\]
In a fixed band about \(y_g\), the damping is \(\ge k_0>0\),
\(|S|\le C\xi(R)\), \(\xi=aR+\tau/R+\tau/R^2\).
Choose the raw radius first, then \(L\) and the coefficient bounds.
The actual shots enter the inner band, and the two boundary signs give
its strict continuation to \(\delta_-/a\). Variation of constants yields
\[
 \int_{R_c}^{R_{\rm out}}|y-y_g|\frac{dR}{R}
 \le C\left(|y(R_c)-y_g|+\delta_-+\frac\tau{R_c}
                                      +\frac\tau{2R_c^2}\right).
\]
The exact first density is
\(Rf_Z=(Bm_R+2\zeta_Rm_Ry+\zeta_RQy^2)/D_y^2\).
At \((3g,1+2g,Q_g,0,y_g)\) it equals \(\nu\), and its derivative
in \(y\) is \(O(g)\). Hence
\[
 |Rf_Z-\nu|\le C(g|y-y_g|+\xi),\qquad
 j(R)=(R/R_c)^\nu U(R),\qquad c\le U,U^{-1}\le C.
\]
Here \(j\) is the linear propagator. Extracting its exact exponent makes the residual density integrable
uniformly along the long corridor, without a \(g^2\log q\) loss.
Rational division also gives
\[
 |f_{ZZ}|\le C(aR^{-2}+gR^{-3}+\tau R^{-4}+\tau R^{-5}),\quad
 |f_{ZZZ}|\le C(aR^{-4}+gR^{-5}+\tau R^{-6}+\tau R^{-7}).
\]
For \(1/3\le\theta=T/T_*<1\), use the coefficient-only radius
\(T_*\), \(r=R/T_*\), \(W=Z/T_*^2\). The exact field is
\[
 r'=g\widehat P(r)+Q(r)W,\quad W'=Wn(r,W),\quad
 \widehat P=(1-r)[a_0+(a_0+a_1)r+w_*r^2],
\]
where \(a_0=q^2\lambda/w_*^2\), \(a_1=q/w_*\),
\(Q=\ell(g-gw_*r)\),
\(n=r(1+2g-gw_*r)-Cg^2w_*(1-w_*r)W\).
On the actual compact strip \(1/3\le r\le1+C\sigma\), \(n\) is a
positive unit. The selected strong graph \(r=\phi(W)\) at \((1,0)\)
has coefficient denominators
\(jB_*+g\widehat c\ge cj\),
\(B_*=1+2g-gw_*\), \(\widehat c=w_*+a_1+2a_0\).
Its power-series majorant gives a fixed normal interval and uniform
\(C^4\) bounds. Put \(x_-=1-\theta\),
\(W_0=g\ell_0x_-\widehat C(\theta)\), and \(u=\phi(W)-r\).
Then \(u_0=x_-(1+O(g))\), with three weighted rows, and exactly
\[
 \begin{gathered}
 (\log u)_{\log W}=-\alpha_--K(W,u),\quad
 \alpha_-=g\widehat c/B_* =\tau/r_-,\\
 |(W\partial_W)^i(u\partial_u)^jK|\le C(W+gu),\quad i+j\le3.
 \end{gathered}
\]
Taylor division of the rational normal field also gives
\(\alpha_-+K\ge c(g+W)\), so the actual orbit reaches \(W=\sigma\)
once. Its integrated densities are
\(\int W\,d\log W\le\sigma\),
\(\int gu\,d\log W\le Cg(1+|\log g|)\).
Use \eqref{deg:scalar-moving-bell} with \(s=\log W\),
\(z=\log u\), \(f=-\alpha_--K(W,e^z)\), and
\(D=E_*:=\theta(1-\theta)\partial_\theta\) at fixed \(s\).
Its initial recursion retains the moving clock
\(s_0=\log W_0\). The two integrated densities above give
\[
 u_\sigma=g^{\alpha_-}x_-^{1+\alpha_-}e^{A(\theta)},\qquad
 |A|+\sum_{j=1}^3|E_*^jA|\le C[\sigma+g(1+|\log g|)],
\]
and \(r_\sigma'=g^{\alpha_-}x_-^{\alpha_-}e^A
(1+\alpha_--x_-A_\theta)>0\).
The fixed normal continuation to \(r=L\) has height
\(H_g(r_0)=\sigma+(L^2-r_0^2)/2+O_{C^3}(g)\), with
\(-H_g'\) a unit. At the coefficient-only cut \(R_c=LT_*\),
its full initial Bell rows are \(T_*^2H_g^{(j)}\).
The same corridor and rational densities give
\(\int|f_{ZZ}U_1|\le Cg\) and
\(\int(|f_{ZZ}U_2|+|f_{ZZZ}U_1^2|)\le Cg\).
Thus the two Bell ratios are bounded, including the third observer row.
For \(\Psi_*(r_0)=\psi(a^2Z_{\rm out}(r_0))\), the exact gain is
\[
 -\Psi_*'=a^2\psi'(a^2Z_{\rm out})T_*^2[-H_g']
                      (\delta_-/(aLT_*))^\nu U(R_{\rm out}),\qquad
 \Omega_*\asymp(aT_*)^{2-\nu}\delta_-^\nu L^{-\nu}.
\]
After removing this exact gain, \(G=(\Psi_*-d_*)/\Omega_*\) and its
inverse have ordinary \(C^3\) bounds, with \(-G'\) a unit. Hence
\[
 \Gamma_-'(b_-)=\frac{\Omega_*}{T_*^2}
 \frac{-G'(r_\sigma)}\theta g^{\alpha_-}x_-^{\alpha_-}e^A
                         (1+\alpha_--x_-A_\theta),\qquad
 |E_*\log\Gamma_-'|+|E_*^2\log\Gamma_-'|\le C.
\]
Neither \(\Omega_*\) nor \(x_-^{\alpha_-}\) has been replaced by a
unit. Factor \(P_-=(1-T/T_*)C_*(T)\), where
\(C_*=1+(1+T_*^{-1})T+\chi T_*T^2\).
On this range \(\mathcal D_-=[C_*/(T^2B_-)]E_*\), with
\(C_*/(T^2B_-)\le C\lambda\). Combining both ranges proves the actual
consumed estimate
\begin{equation}\label{deg:far-negative-budget}
 |L_-|\le C\{\tau(1+|\log\tau|)^3+\lambda\},
                       \qquad T_b\le T<T_*.
\end{equation}
In the root range, the residual logarithmic derivatives are bounded,
and the factor \(C_*/(T^2B_-)\) makes their action derivatives small.
Together with the positive-section and relative-variation estimates,
this gives
\begin{equation}\label{deg:far-bell}
 |\varrho|+\max_{i+j\le2}|\mathcal D_+^i\mathcal D_-^j\varrho|
 +|L_-|\le C\{\tau(1+|\log\tau|)^3+\lambda\},\qquad
 |L_+|\le C_0\sigma+C_\sigma[X_0+\tau(1+|\log\tau|)^3].
\end{equation}
Therefore the full physical identity is
\begin{equation}\label{deg:far-curvature}
 \frac{UTP_+\phi_+'}{U+T}\frac{\widehat\Pi''}{\widehat\Pi'}
 =3\chi UT-1+E_f,\quad
 |E_f|\le\epsilon_0+C_\sigma[\tau(1+|\log\tau|)^3+\lambda],
 \quad\epsilon_0<1/8.
\end{equation}
Independent \(q_0,\lambda_0,\tau_0\) make this strictly negative.
The exact endpoint powers also remain: \(x_-=T_*-T\) obeys
\(x_-\sim C_\theta(A_+-U)^{(r_++\tau)/(r_-+\tau)}\), while
the two physical embeddings carry powers \(1+\tau/r_\pm\).
The physical exponent is \(r_+/r_-\asymp3\lambda\); only weighted
endpoint rows and pointwise \(C^3\) on the open passage are used.

It remains to establish raw admission and the whole upper complement.
The two comparison shots \(U=A_+-X_j\) have positive-unit heights at
fixed cuts \(L_+>A_++1\), \(L_->\max T_j+1\), since their fast
invariants give \(Z_+(L_+)=\tfrac12[L_+^2-(A_+-X_j)^2]+O(\tau(1+|\log\tau|))\)
and \(Z_-(L_-)=\tfrac12(L_-^2-T_j^2)+O(\tau(1+|\log\tau|))\).
Choose a fixed upper fast shot and \(\eta_{\max}>0\); scalar order
bounds every intervening central shot through that height at both cuts.
Before applying an observer, compare their ratios to the passive
parabolic references
\[
 t(Y_{{\rm par},\pm})_t=Y_{{\rm par},\pm}
 \left\{\frac{1\pm t+Ct^2Y_{{\rm par},\pm}}
 {\pm t+\ell(\pm t)Y_{{\rm par},\pm}}-2\right\},\qquad
 Y_{{\rm par},\pm}(0)=1/2.
\]
Their analytic recursion has denominator \(j+2\), and the original
prepared passive compact set retains their strict raw-window margins.
The actual/reference ratio fields have negative \(y\)-derivative and
difference \(O(g+s^2/t^2+sp/t)\). Logarithmic variation of constants
at each fixed raw radius yields
\[
 |y(\delta_\pm)-Y_{{\rm par},\pm}(\delta_\pm)|
 \le C[(aR_c/\delta_\pm)^{c_0}+g+\tau/R_c+\tau/R_c^2].
\]
Thus all comparison and upper core shots are admitted uniformly in the
independent rectangle. For the entire lower output family, the raw-height
spectral slope \eqref{deg:spectral-exterior-slope}, with identity observer, gives
\(|h_T|\le Ca^{2-\nu}\delta_-^\nu T^{1-\nu}\), hence total variation
\(\le C(aT_*)^{2-\nu}\delta_-^\nu\). The raw version of the root-range
gain gives the same bound. Since \(aT_*\le4g\), the total is
\(O(g^{3/2})\), smaller than the fixed reference buffer. The bounded
positive normal family has gain \(O(a^{2-\nu})\) and the same input
admission. Ordinary observers are composed only on these admitted images;
the original smaller windows, speed and stop clips are then intersected.

On the fixed core, \(D=\tau P+QZ\) and
\(f_Z=d\log(Z/D)/dU+[\tau P'+(Q'+\zeta)Z]/D\).
The inner positive piece has \(P\ge p_c\), the outer positive piece
has the comparison's \(Z\ge z_c\), and the negative fixed slab has
\(P_-\ge p_c\). The quotient is uniformly integrable and endpoint
\(Z/D\) are units, giving bounded core variation and inverse.
The positive tail has \(m_+=\tau/R^2-\tau/R+3g+aR\),
\(B_+=1+2g+aR\), \(Q_+=\ell(g+aR)\),
\(\zeta_{R,+}=CaR(g+aR)\); the same exact \(y_g,\nu\) corridor
applies. Both absolute outward gains are
\(a^{2-\nu}\delta_\pm^\nu\) times units. Their common factor cancels
in the actual \(\Pi'\), so \(0<\Pi'\le C_c\) through
\(\eta_{\max}\). All ordinary observer and divided closing rows are
those of Lemma~\ref{lem:closing}; their third chain rule uses only
\(f''',d'''\), retaining each small gain.

The action gap gives
\(0<\eta_{c,1}<\eta_{c,0}\le C\tau^{-N}e^{-c/\tau}\).
Fixing the shots and caps first, the independent shrink gives
\begin{equation}\label{deg:far-physical-cover}
 \begin{gathered}
 \eta_{c,1}<\eta_{c,0}<\eta_H:=\tau K_H(1+\tau)<\eta_{\max}/2,\\
 H_H=a^2\eta_H=K_H(s^2+(sp)^2),\qquad a^2\eta_{\max}<H_{\rm cap}/2.
 \end{gathered}
\end{equation}
On the original height axis, lower through the \(X_0\) shot, core
starting at the \(X_1\) shot, and high above \(H_H\) cover \(I_\theta\)
with strict overlaps. The positive equality \(H=H_H\) belongs to the
core; scalar order retains all clipped nonroot inputs. Choose \(m\)
last so \(m\max(C_c,C_H)<1/2\). The upper two packages have
\(F'>1/2\); on the lower one \(\Pi''<0\), so every critical point
of the same \(F\) has \(F''>0\). Lemma~\ref{sharp:upper-critical-count}
proves the whole-fibre bound, excluding its natural stopped endpoint.
\end{proof}

\subsection{Cubic and quadratic source compactifications}
\label{subsec:deg-cubic-quadratic}

The chart \eqref{deg:cubic-chart} covers the compact cubic limits of
the small-\(p\) region. Proposition~\ref{deg:prop-compact-cubic}
proves their complete return and whole-fibre count. The estimates below
supply its growing tails, where a compact phase comparison no longer
applies. We then return to the bounded positive quadratic face of
\eqref{deg:gamma-chart}.

On a growing positive tail, the lip label is unbounded while its raw
abscissa remains small. The next estimate controls the logarithmic
derivatives of the actual exterior map in that regime. Keeping its
absolute first gain will later allow composition with the independent
negative section map.

\begin{lemma}[Complete section derivatives on a growing cubic tail]
\label{deg:lem-raw-growing-section}
Fix a sufficiently small terminal cut \(\delta>0\). Let
\[
 \begin{gathered}
 M(x)=x^3+\mu_2x^2+\mu_1x+\mu_0,\qquad B_0(x)=1+b+x,\\
 Q_0(x)=1-A(v+x)+C(v+x)^2,\qquad \zeta_0(x)=C(v+x).
 \end{gathered}
\]
with bounded \(A,C\). Hold all coefficients fixed as the starting
point \(a\) varies. On an enlarged interval \(I\subset(0,d_0]\),
with \(d_0<\delta/2\), suppose uniformly for \(a\in I\)
\[
 |b|+|v|\le C_0a,\qquad |\mu_j|\le C_0a^{3-j}\ (j=0,1,2),
 \qquad M(x)\ge c_0x^3\quad(a\le x\le\delta).
\]
Use the actual positive graph
\begin{equation}
 z_x=\frac{z\{xB_0(x)+\zeta_0(x)z\}}{D_0(x,z)},\qquad
 D_0=M+Q_0z,\qquad z(a)=\ell_0M(a),
 \label{deg:weighted-raw-positive-graph}
\end{equation}
and an increasing regular terminal adapter \(\chi\), with uniformly bounded
\(C^3\) norm and inverse on its enlarged image. For fixed positive
\(\sigma\), label the input by \(b_+=-a^2/(2\sigma^2)\), and write
\(\Gamma_+(b_+)=\chi(z_\delta)\). After choosing \(\ell_0,d_0\)
small, this continuation and every separate affine normalization
\(\phi_+\) satisfy
\begin{equation}
 |(a\partial_a)^j\log\phi_+'|\le Ca(1+|\log a|)^4,
 \qquad j=1,2.
 \label{deg:raw-growing-section-jets}
\end{equation}
The exact positive factor \(\sigma^2\) in the first derivative is
retained; no bound on its reciprocal is required.
\end{lemma}
\begin{proof}
Set \(E=a\partial_a\) and \(m(x)=M(x)/x^2\). On \(a\le x\le2a\),
the coefficient bounds and the polynomial margin give
\[
 c_0a\le m(x)\le C_1a,\qquad
 |b|+|v|+|\mu_2|+a+|\mu_1|/a+|\mu_0|/a^2\le C_1a.
\]
Thus Lemma~\ref{deg:lem-spectral-exterior} applies with \(t=1\)
and comparison mass \(\rho=a\), including its near-lip strip entry.
This mass is not differentiated;
all polynomial coefficients and \(\sigma\) remain fixed. The adapter
is evaluated on its stated enlargement of the generated terminal image.
In the actual raw-section applications,
\eqref{deg:spectral-exterior-window}, with \(\rho=a\), and the
original fixed terminal-image margin admit the whole family after
choosing \(d_0\) within its original cutoff reserve. The stated
ordinary \(C^3\) adapter and inverse bounds give a uniform bound for
its two logarithmic-height rows on the enlarged terminal image.

The common spectral estimate has
\[
 \lambda=\frac{4\mu_2}{1+b}=O(a),\qquad
 \Theta=a+|\mu_1|/a+|\mu_0|/a^2\le C_1a.
\]
Its terminal coefficient \(\mathcal T\) is bounded on the stated
buffer, and \(\mathcal T(a/\delta)^{2-\lambda}\le Ca\), after
the permitted \(d_0\) shrink. Consequently
\[
 |E\log\Gamma_+'+\lambda|+|E^2\log\Gamma_+'|
 \le Ca(1+|\log a|).
\]
The physical gain is given exactly by
\eqref{deg:scalar-polynomial-lip-slope},
\begin{equation}\label{deg:weighted-complete-positive-derivative}
 \Gamma_+'=\sigma^2\chi'(z_\delta)\frac{z_\delta}{D_0(\delta,z_\delta)}
                 B_0(a)J_0(a)e^{I(a)},\qquad J_0=1+O(a),
\end{equation}
where \(J_0=K\) and \(I\) are the exact factors in that formula.
This retains the factor \(\sigma^2\) without a reciprocal bound.
Since \(\lambda=O(a)\), both uncorrected logarithmic rows satisfy
the stronger first-logarithmic-power bound. Frozen affine calibration
leaves them unchanged, proving \eqref{deg:raw-growing-section-jets}.
\end{proof}

\begin{lemma}[Weighted cubic tail on the actual sections]
\label{deg:lem-weighted-cubic-tail}
Fix a compact set of coefficients \(0\le U\le U_0\),
\(|G|\le G_0\), and use the field
\eqref{deg:compact-cubic-field}--\eqref{deg:compact-cubic-polynomials}.
On its root-free positive fibres choose \(M\) sufficiently large,
independently of \(\tauc\), and retain
\[
 u\ge2M,\qquad a=\tauc u\le d_0,\qquad r_0\le r<c.
\]
Here \(c\) is the simple negative root and the two moving sections are
\(Y=\ell_0P\), with \(\eta=\tauc Y\).  The exterior maps are the
actual continuations to the fixed raw sections, followed by the regular
input adapter and the divided finite-side output adapter, respectively.
These continuations use the nested domains constructed by
\eqref{deg:compact-raw-corridor}.
There are constants, uniform on these domains, such that, with
\[
 E=u\partial_u,\quad
 \mathcal D_+=\frac{P(u)}{uB(u)}\partial_u,\quad
 \mathcal D_-=\frac{P(-r)}{rB(-r)}\partial_r,\quad
 \epsilon_u=\tauc(1+u),\quad L_u=1+|\log\epsilon_u|,
\]
the exact action matching has the form
\begin{equation}
 A_+(u)-A_-(r)=\tauc\{\log P(u)-\log P(-r)+R_+(u)-R_-(r)\},
 \label{deg:weighted-split-action}
\end{equation}
where
\begin{equation}
 |R_+|+|R_-|\le C,\qquad
 |E^jR_+|\le C\epsilon_u,\qquad
 |\mathcal D_-^jR_-|\le C\tauc\quad(1\le j\le3).
 \label{deg:weighted-action-jets}
\end{equation}
The relative variation in \eqref{deg:weighted-relative-variation} obeys
\begin{equation}
 |\varrho|+\frac{|\mathcal D_+\varrho|+|\mathcal D_-\varrho|}{u+r}
       \le C\epsilon_uL_u^4.
 \label{deg:compact-rho}
\end{equation}
If \(b_+=-u^2/2\),
\(b_-=-r^2/2\), and \(\phi_\pm\) are the separately affine-normalized
actual section embeddings, then
\begin{equation}
 |E\log\phi_+'|+|E^2\log\phi_+'|
       \le C\epsilon_uL_u^4,
 \qquad
 |\mathcal D_-^j\log\phi_-'|
       \le C\tauc(1+|\log\tauc|)^4\quad(j=1,2).
 \label{deg:weighted-actual-section-jets}
\end{equation}
Every prime on \(\phi_\pm\) is taken in its own label \(b_\pm\).
The full composition \(\Pi=\phi_-\circ S\circ\phi_+^{-1}\), with
\(S(-u^2/2)=-r^2/2\), therefore satisfies
\begin{equation}
 \phi_+'\frac{P(u)}{u+r}\frac{\Pi''}{\Pi'}
       =3-\frac{U}{ur}+E_u,
 \qquad |E_u|\le C\epsilon_uL_u^4.
 \label{deg:compact-curvature}
\end{equation}
The inward action estimates and the two exterior estimates are obtained
separately on their actual domains.
\end{lemma}

\begin{proof}
We first continue the two inward graphs to the original centre and
differentiate their exact height match. We then estimate the two
exterior embeddings separately. The negative embedding requires a
root-collar calculation because its first variation vanishes there.
Substitution in the completed-curvature identity gives the final estimate.

Choose a small fixed \(\delta_c>0\) such that \(P\) is a positive unit
on \([-2\delta_c,2\delta_c]\), and so that
\(2\delta_c<r_0\).  On the root-free positive fibre, put
\[
 \begin{gathered}
 q=\frac{A_+(u)-A_+(v)}\tauc,\quad
 W=\frac{Y}{P(v)},\quad D(v)=\frac{P(v)}{vB(v)},\\
 H(v)=\frac{P'(v)}{vB(v)},\qquad
 \varpi(v)=\frac{C\tauc^3(G+v)P(v)}{vB(v)}.
 \end{gathered}
\]
For \(\delta_c\le v\le u\), graph division gives the exact equations
\begin{equation}
 v_q=-\tauc D(v),\qquad
 W_q=W\left\{-\frac{1+\varpi W}{1+QW}+\tauc H\right\},
 \qquad W(0)=\ell_0.
 \label{deg:weighted-inward-equations}
\end{equation}
In particular, \(\tauc H=O(\epsilon_u)\),
\(Q-1=O(\epsilon_u)\), and
\(\varpi=O(\epsilon_u^3)\).  Fix \(\ell_0\), then \(d_0\),
small enough that the braces and the derivative of the right-hand side
with respect to \(W\), for \(0\le W\le\ell_0\), are both less than
\(-c_1<0\).  A first-exit argument gives
\begin{equation}
 0<W(q)\le\ell_0e^{-c_1q}.
 \label{deg:weighted-kernel-decay}
\end{equation}

For unbounded \(u\), Euler differentiation keeps the moving entrance
under control. At fixed action clock \(q\),
\begin{equation}
 Ev=\kappa(u)D(v),\qquad
 \kappa(u)=\frac{u^2B(u)}{P(u)}\asymp u^{-1}.
 \label{deg:weighted-transport}
\end{equation}
The first three Euler derivatives of \(\kappa\) are \(O(u^{-1})\).
On \([M,u]\), \(D(v)=O(v^2)\), and each further application of
\(D(v)\partial_v\) increases polynomial degree by at most one.
On \([\delta_c,M]\), this vector field and its first four derivatives
are uniformly bounded, including when the positive factor approaches a
double root: they contain \(P\), rather than \(P^{-1}\).
Differentiating \eqref{deg:weighted-transport} inductively therefore gives
\[
 |E^j\log v|\le C_j\quad(1\le j\le4),\qquad
 |E^j(\tauc H)|+|E^j(Q-1)|+|E^j\varpi|
       \le C_j\epsilon_u\quad(0\le j\le3).
\]
Because the transported field contains \(P\), these bounds remain
uniform when the root-free positive minimum tends to zero.
Apply Lemma~\ref{deg:lem-stable-ratio-ledger} with
\(D=E,\kappa=\tauc,\rho=\epsilon_u,\alpha=0\) on
\(\delta_c\le v\le u\).
The needed clock rows follow from \(\partial_q=-\tauc D(v)\partial_v\):
on the growing part its degree growth is controlled by
\(\tauc v\le\epsilon_u\), and on the compact part its coefficients
are bounded. The strict diagonal chosen above gives
\begin{equation}\label{deg:weighted-kernel-jets}
 |E^jW(q)|\le C\epsilon_u(1+q)^Ne^{-c_2q},\qquad1\le j\le3.
\end{equation}

The nonlinear logarithmic correction on this tail is exactly
\[
 V_+(u)=\int_0^{q_c(u)}
       \frac{(Q-\varpi)W}{1+QW}\,dq,
 \qquad q_c(u)=\frac{A_+(u)-A_+(\delta_c)}\tauc.
\]
The fixed-clock and terminal formulas of the same lemma give \(|V_+|\le C\) and
\(|E^jV_+|\le C\epsilon_u\), \(1\le j\le3\): indeed,
\(|E^jq_c|\le C_j/\epsilon_u\) and
\(q_c\ge[A_+(2M)-A_+(M)]/\tauc\ge a_M/\tauc\).
Every boundary term is therefore at most
\(C\tauc^{-N}e^{-c_2a_M/\tauc}\).

For the negative inward tail replace \((P,B,Q,\varpi)\) by
\((P(-v),B(-v),Q(-v),-C\tauc^3(G-v)P(-v)/(vB(-v)))\).
At fixed \(q\), \(\mathcal D_-\) transports to
\(P(-v)[vB(-v)]^{-1}\partial_v\);
its coefficient rows in the braces are \(O(\tauc)\) on
\([\delta_c,c]\), including the root factor. The same lemma with
\(\rho=\tauc\) gives \(|V_-|\le C\),
\(|\mathcal D_-^jV_-|\le C\tauc\). Since \(r\ge r_0>2\delta_c\),
its complete terminal rows are exponentially small.

Continue both tails to \(\xi=0\) using
Lemma~\ref{deg:lem-compact-height-continuation}, with radial coordinate
\(v=|\xi|\) oriented from \(\delta_c\) to zero and
\(\kappa_c=\tauc\). Here \(P\ge p_c>0\), the linear inward
propagator is at most one, and the nonlinear coefficient has ordinary
height derivatives through order three bounded by \(C/\tauc\).
For the incoming heights \(y_\pm=Y(\pm\delta_c)\), the exact memory
\eqref{deg:stable-ratio-polynomial-memory} gives
\(y_\pm\le C\tauc^{-N}e^{-c/\tauc}\); its three
logarithmic phase rows grow at most polynomially in \(\tauc^{-1}\).
The cited continuation and composition rows therefore make the remaining
nonlinear correction exponentially small, with those phase derivatives.
The exact cocycle \eqref{deg:exact-action-cocycle} gives the actual
heights on the common section:
\begin{align*}
 \log Y_0^+(u)&=\log(\ell_0P(u))-A_+(u)/\tauc+R_+(u),\\
 \log Y_0^-(r)&=\log(\ell_0P(-r))-A_-(r)/\tauc+R_-(r).
\end{align*}
with \eqref{deg:weighted-action-jets}.  Equality of these heights is
exactly the physical matching condition.  This proves
\eqref{deg:weighted-split-action}; it also supplies a separated extension
\(R_0(u,r)=R_+(u)-R_-(r)\), whose genuinely mixed derivatives vanish.

The separated height formulas also determine the relative first
variation and its weighted derivatives. Write
\(P_+=P(u)\), \(P_-=P(-r)\), \(B_+=B(u)\), \(B_-=B(-r)\),
and use radial derivatives in
\(H_+=P_+'/(uB_+)\), \(H_-=P_-'/(rB_-)\).
Differentiating \eqref{deg:weighted-split-action} gives
\begin{equation}
 r_u=\frac{uB_+P_-}{rB_-P_+}(1+\varrho),\qquad
 1+\varrho=
 \frac{1-\tauc H_+-\tauc\mathcal D_+R_+}
      {1-\tauc H_--\tauc\mathcal D_-R_-}.
 \label{deg:weighted-relative-variation}
\end{equation}
Since \(\mathcal D_+=(P_+/(u^2B_+))E\), the numerator differs
from one by \(O(\epsilon_u)\), with the same Euler bound.  The
denominator differs from one by \(O(\tauc)\), with two bounded action
variations of that order.  This proves \eqref{deg:compact-rho} (even
with \(L_u^4\) omitted).  Both denominators are positive after shrinking
\(d_0\) and then \(\tauc\).

For the positive exterior map apply
Lemma~\ref{deg:lem-raw-growing-section} with
\[
 \sigma=\tauc,\quad a=\tauc u,\quad
 (\mu_2,\mu_1,\mu_0)=(3\tauc G,-\tauc^2U,\tauc^3),
 \quad b=2\tauc G,\quad v=\tauc G.
\]
For \(u\ge2M\), the coefficient bounds hold and the leading cubic
term gives \(M(x)\ge c x^3\) on \([a,\delta]\), uniformly on the
fixed compact block. In the translated coordinate \(x=v-\gamma\),
the original positive raw port is \(x_T=\delta-\tauc G\).
The proof of the raw lemma is uniform for terminal cuts in a fixed
neighborhood of \(\delta\); \(x_T\) is coefficient-only and is held
fixed under \(E\). Thus it applies at this actual port. The corridor
supplies its enlarged terminal image.
Since \(a\asymp\epsilon_u\), \eqref{deg:raw-growing-section-jets}
is the positive assertion of \eqref{deg:weighted-actual-section-jets}.

Put \(\lambda_G=12\tauc G/(1+2\tauc G)\).
For the negative exterior map, first restrict to
\(r_0\le r\le c-x_0\), with \(x_0>0\) fixed and small.
Here \(P_-(r)=P(-r)\) has a fixed positive margin. Put \(v=-\xi\),
\(B_-(v)=B(-v)\), \(Q_-(v)=Q(-v)\), and use the increasing
height clock \(s=\log\eta\). The fast invariant
\(H=\eta-v^2/2\) satisfies the exact equation
\[
 H_s=\tauc\mathscr F(H,e^s),\qquad
 \mathscr F=
 \frac{v\eta(B_--Q_-)/\tauc
       -C\tauc(G-v)\eta^2-vP_-(v)}
      {vB_-(v)-C\tauc^2(G-v)\eta}.
\]
On a fixed box with \(v\ge r_0/2\), the denominator is a positive
unit, \((B_--Q_-)/\tauc\) is smooth and bounded, and
\(v=\sqrt{2(e^s-H)}\) gives
\( |\partial_H^i\partial_s^j\mathscr F(H,e^s)|\le C\)
for \(i+j\le3\). The initial rows are
\[
 s_l=\log\{\tauc\ell_0P_-(r)\},\qquad
 H_l=-r^2/2+\tauc\ell_0P_-(r).
\]
The ordinary \(r\)-derivatives of \(s_l-\log\tauc\) are bounded.
The moving initial recursion of
\eqref{deg:scalar-moving-bell} therefore gives the three initial
variations of \(-r^2/2\), with errors \(O(\tauc)\).
Choose a fixed small \(\eta_*>0\). First exit and the three
triangular equations on
\(s_l\le s\le\log\eta_*\), whose length is
\(O(1+|\log\tauc|)\), give
\[
 H(\log\eta_*;r)=-r^2/2+
          O_{C^3_r}\bigl(\tauc(1+|\log\tauc|)\bigr).
\]
This is an actual outward passage: initially
\(D=\tauc P_-+Q_-\eta>0\), and at a possible first zero
\(D_s=Q_-\eta>0\). On the fixed positive-height continuation
to a sufficiently large fixed cut \(v=L>c+1\), the graph equation
is an \(O_{C^3}(\tauc)\) perturbation of \(\eta_v=v\). Hence
\[
 \eta_L(r)=\tfrac12(L^2-r^2)+
          O_{C^3_r}\bigl(\tauc(1+|\log\tauc|)\bigr).
\]
The cut and its buffered chord lie in the original corridor
\eqref{deg:compact-raw-corridor}. From this cut to the raw port use
the negative tail with source quartet
\((a,p,\gamma,s)=(-\tauc,\tauc U,\tauc G,\tauc)\).
For a fixed \(0<\nu<1\), \eqref{eq:N-bound} gives
\[
 \int_L^{\delta/\tauc+G}\|\mathcal N\|_{C^3}\,dv
 \le C\tauc\int_L^\infty
 (v^{-2+\nu}+v^{-3+\nu}+v^{-4+\nu}+v^{-5+\nu})\,dv
 \le C\tauc.
\]
After extraction of its exact first slope, the tail and its inverse
are \(\mathrm{id}+O_{C^3}(\tauc)\). Its physical slope is
\(d_{\rm c}=\tauc^{2-\lambda_G}U_{\rm c}\), with
\(U_{\rm c}\) a positive unit. The original ordinary \(C^3\)
observer and divided closing add \(O_{C^3}(d_{\rm c})\) by
\eqref{eq:exterior-normalization}. Calibrating value and slope at a
frozen reference in this compact interval gives
\(\widetilde\phi_-=\mathrm{id}+O_{C^3}(\tauc(1+|\log\tauc|))\)
in \(b_-=-r^2/2\). Since \(\mathcal D_-\) has bounded coefficients
and derivatives here, the two logarithmic slope rows have the same
bound. They are unchanged by a frozen positive affine calibration,
so the original \(\phi_-\) satisfies
\eqref{deg:weighted-actual-section-jets} on this interval.
Near the root, the first variation vanishes and the preceding compact
normalization no longer gives a bounded inverse. We extract its endpoint
power before estimating logarithmic derivatives.

The root equation gives the exact factorization
\[
 P_-(r)=(c-r)L_-(r),\qquad
 L_-(r)=r^2+(U/c+c^{-2})r+c^{-1},\qquad
 L_-(c)=c^2+U+2/c>0.
\]
Thus \(c\) and \(L_-\), with their required finite coefficient
derivatives, are uniform units on a fixed negative collar,
independently of the positive bottleneck. Put \(x=c-r\), \(y=\eta\).
Dividing by the
positive normal rate \(rB(-r)\) gives
\[
 \dot x=-\beta x p(x)-k_1(x)y,\qquad
 \dot y=y\{1+k_2(x)y\},\qquad
 \beta=\frac{\tauc L_-(c)}{cB(-c)}=O(\tauc),
\]
Here the coefficients are exactly
\[
 \begin{split}
 p(x)&=\frac{cB(-c)L_-(c-x)}{L_-(c)(c-x)B(-c+x)},\\
 k_1(x)&=\frac{Q(-c+x)}{(c-x)B(-c+x)},\qquad
 k_2(x)=-\frac{C\tauc^2(G-c+x)}{(c-x)B(-c+x)}.
 \end{split}
\]
In particular \(p(0)=1\), \(p,k_1\) are positive units and
\(k_2=O_{C^7}(\tauc^2)\) on a fixed box.
Use the normal-time construction
\eqref{eq:saddle-normal-time}--\eqref{eq:exact-saddle-invariant} in the
proof of Lemma~\ref{lem:actual-dulac-matching}, with
\(\rho=\beta\), \(h=-\beta xp\), \(c=-k_1\), \(k=k_2\),
and selected coefficient parameter \(\tauc\) at fixed \(U,G,A,C\).
The explicit coefficients have a bounded joint order-seven budget:
\(h=O_{C_x^7}(\beta)\), \(k_1\) is a positive unit, and their
selected \(\tauc\)-rows are bounded. The normal coefficient and its
parameter rows retain their height factors.
The projection at phase order five uses axis-flow order six and the
selected parameter row. For a fixed normal decay rate \(0<\lambda<1\),
its integrands have the majorant
\(C(1+T)^Ne^{-(\lambda-C\beta)T}\), integrable after
\(C\beta<\lambda/2\). Thus the exact coordinates \((\Xi,Z)\)
and invariant \(I=\Xi Z^\beta\) have five phase derivatives and one
intrinsic \(\tauc\)-derivative, as required for the divisions below.
In the following coefficient-dependent units, a subscript
\({\rm c}\) records dependence on \(\tauc\). On the moving curve
\(y=\tauc\ell_0P_-\), these coordinates have the form
\[
 \Xi_m=xu_{\rm c}(x),\qquad Z_m=\tauc xv_{\rm c}(x),\qquad
 I(x)=\theta_{\rm c} x^{1+\beta}A_{\rm c}(x),\quad
 \theta_{\rm c}=\tauc^\beta,\quad A_{\rm c}=u_{\rm c} v_{\rm c}^\beta,
\]
with positive \(C^4\) units \(u_{\rm c},v_{\rm c},A_{\rm c}\).
To check the coefficient differences before taking the root limit, put
\(Y_{\rm c}(x)=\tauc x\ell_0L_-(c-x)\). Then
\[
 \partial_{\tauc}u_{\rm c}(x)=
 \frac{(\partial_{\tauc}\Xi)(x,Y_{\rm c}(x))}{x}
 +\Xi_y(x,Y_{\rm c}(x))\ell_0L_-(c-x),\qquad
 v_{\rm c}=\ell_0L_-(c-x)e^{-J(x,Y_{\rm c}(x))}.
\]
The first numerator vanishes at \(x=0\); its parameter-valued fifth
phase row gives \(C^4\) after Taylor division. In the second term the
moving factor \(x\) has already cancelled, so \(\Xi_y\) needs only
four phase rows. The formula for \(v_{\rm c}\) removes the factor
\(\tauc x\) before differentiation. Integrating these intrinsic
parameter rows gives \(u_{\rm c}-u_0,v_{\rm c}-v_0=O_{C^4}(\tauc)\),
and hence \(A_{\rm c}-A_0=O_{C^4}(\tauc)\). Also
\(\theta_{\rm c}=1+O(\tauc|\log\tauc|)\).

At a fixed outer cut beyond \(-c\), the stopped trajectory has
positive height. The invariant is a regular increasing coordinate
there, including \(I=0\), as in
\eqref{eq:exact-saddle-invariant}. Choose its buffered invariant chord
first, then the smaller lip collar so \(I(x)\) lies strictly inside it.
Choose the upper horizontal boundary with invariant values above that
chord throughout the local strip. This is possible at \(\tauc=0\),
where \(\pi=c-\sqrt{r^2-2\eta}\), and the strict margin persists
for small \(\tauc\). Thus no orbit of the lip collar exits through
that boundary.
The outward speed starts positive. At a possible first zero of
\(\mathscr D=\tauc P_-+Q_-\eta\), one has
\(\dot r=0\) and \(\dot{\mathscr D}=Q_-\dot\eta>0\), excluding
loss of that speed. Increasing height and the invariant chord then give
the actual first hit of the outer cut at positive height, uniformly for
\(0<x\le x_0\). Continue this fixed chord to the
negative raw section, before either the regular output adapter or the
finite closing map. After extracting the intrinsic exact first slope, write
\[
 T(x)=t_{\rm st}+d_{\rm c} F_{\rm c}(I(x)),\qquad
 d_{\rm c}=\tauc^{2-\lambda_G}U_{\rm c}>0,
\]
where \(U_{\rm c}\) is a positive unit and \(F_{\rm c},F_{\rm c}^{-1}\)
have uniform \(C^4\) bounds on a fixed invariant chord. Here the
finite positive-height passage has \(O(\tauc)\) coefficient changes.
On its subsequent tail, \eqref{eq:N-bound}, integrated with the
frozen negative quartet, bounds the nonlinear normalized kernel by
\(O(\tauc)\) through four phase orders. Thus the intrinsic,
separately normalized maps
satisfy \(F_{\rm c}=F_0+O_{C^3}(\tauc)\); no convergence rate for
the extracted slope \(d_{\rm c}\) is used.

Since \(b_-=-r^2/2\), one has \(db_-/dx=r=c-x>0\).
Consequently the exact first derivative is
\[
 T_{b_-}'=d_{\rm c}\theta_{\rm c} x^\beta
                 K_{\rm c}(x,x^\beta),\qquad
 K_{\rm c}(x,z)=
 \frac{F_{\rm c}'(\theta_{\rm c} xzA_{\rm c}(x))
       \{(1+\beta)A_{\rm c}(x)+xA_{\rm c}'(x)\}}{c-x}.
\]
The function \(K_{\rm c}\) is a positive \(C^3\) unit on a fixed
rectangle \(0\le x\le x_0\), \(0\le z\le1\), and
\(K_{\rm c}(0,z)\) is independent of \(z\).
At \(\tauc=0\) the outer graph is
\(\eta(\xi)=(\xi^2-r^2)/2\). Its intrinsic, separately
normalized section map is affine in \(b_-\). The regular factors
therefore cancel exactly at \(z=1\), giving a positive constant
\(k_0\), independent of \(x\), for which
\[
 \|K_{\rm c}(\cdot,1)-k_0\|_{C^2}
       \le C\tauc(1+|\log\tauc|).
\]
The affine limiting map is what makes the logarithmic derivative small;
boundedness of the regular factors alone would not suffice.

To differentiate along \(z=x^\beta\), use
\(E_x=x\partial_x+\beta z\partial_z\). For
\(0<x\le x_0<1\),
\[
 x|1-x^\beta|\le C\beta.
\]
The uniform \(C^3\) bounds for \(K_{\rm c}\), its constancy in
\(z\) at \(x=0\), and the preceding cancellation imply
\[
 |E_x\log K_{\rm c}|+|E_x^2\log K_{\rm c}|
       \le C\tauc(1+|\log\tauc|)
       \quad\hbox{on }z=x^\beta.
\]
Indeed the differences of the \(x\)-derivatives between \(z\)
and \(1\) have factors \(1-z\), accompanied by \(x\) or
\(x^2\); all remaining terms contain \(\beta\).
Since
\[
 \mathcal D_-=-\frac{L_-(c-x)}{(c-x)B(-c+x)}E_x,
 \qquad E_x\log x^\beta=\beta,
\]
we obtain the first two action-logarithmic bounds for
\(T_{b_-}'\), including their one-sided limits as \(x\downarrow0\).
The derivative itself may vanish like \(x^\beta\).

Finally append the original increasing ordinary \(C^3\) output adapter
\(\psi\) and the divided closing map. Put
\[
 \Psi_{\theta,m}=\mathfrak f_{\theta,m}\circ\psi,
 \qquad \mathfrak f_{\theta,m}(q)=f_\theta(mq)/m,
 \qquad \Gamma_-=\Psi_{\theta,m}\circ T.
\]
Admission to the original raw and closing windows precedes these
compositions. Apply \eqref{phys:eq-observer-log} and
\eqref{phys:eq-closing-jets} with \(D=\mathcal D_-\) and primes in
\(b_-\). They bound the two logarithmic-slope corrections using only
the ordinary \(C^3\) norms and positive first derivative of
\(\Psi_{\theta,m}\), uniformly for \(0\le m\le m_0\).
The invariant formula retains the small gain:
\(|DT|+|D^2T|\le Cd_{\rm c}\), and
\(d_{\rm c}\le C\tauc^2\), since \(\lambda_G=O(\tauc)\).
Both additional logarithmic rows are therefore \(O(d_{\rm c})\),
using only the ordinary \(C^3\) adapter and closing bounds of
Lemma~\ref{phys:lem-c3-adapters}. Thus the actual complete output
satisfies \eqref{deg:weighted-actual-section-jets}; affine normalization
fixes its logarithmic rows.

We can now combine the inward match with both exterior estimates.
Apply Lemma~\ref{neg:lem-completed-curvature} with
\(v=r\), \(P_+=P(u)\), \(P_-=P(-r)\), and \(W=u+r\).
In \eqref{deg:weighted-exact-curvature-chain}, put \(R_B=B_+/B_-\).
The first term differs from
\((P'(u)/u+P'(-r)/r)/(u+r)=3-U/(ur)\) by
\(O(\epsilon_u)\).  The second is \(O(\epsilon_u)\), since
\(R_B-1=\tauc(u+r)/B_-\) and
\(r_u=O(P_-/(u^2r))\).  The third is controlled by
\eqref{deg:compact-rho}, and the last by
\eqref{deg:weighted-actual-section-jets}, since
\(P_+/(u^2(u+r))\) is bounded.  For the negative section term use
\[
 \frac{\tauc(1+|\log\tauc|)^4}{u+r}
       \le C\epsilon_u(1+|\log\epsilon_u|)^4.
\]
To see this uniformly, write \(\epsilon_u=\tauc(1+u)\), use
\(1+|\log\tauc|\le1+|\log\epsilon_u|+\log(1+u)\), and note that
\((1+\log(1+u))^4/(1+u)^2\) is bounded.
This proves \eqref{deg:compact-curvature} in the actual affine
input coordinate.
\end{proof}


\begin{proposition}[Positive quadratic face]
\label{deg:prop-positive-quadratic}
Fix \(\Pi<\infty\). For \(\gamma=g>0\),
\(0<k<k_0\), and \(0<\pi\le\Pi\), the actual low-height through
interval on the real/fold side and its fixed adjacent root-free band
has the moving-section descriptions below. Every zero of the complete
map's curvature on this interval is simple and is a downward crossing
in the same increasing physical input. The bounded-height complement
has uniformly bounded complete derivative. On the root-free side
outside this band, the complete return is a contraction after a
coefficient shrink.
The domains and estimates persist through the actual discriminant
\(\Delta=0\), whose limit at \(k=0\) is \(\pi^2=12\).
\end{proposition}

\begin{proof}
We use the curvature principal part \(3-\pi/(uR)\), with the
homogeneous section powers retained. Its zeros require separate
arguments in the joint growing corner and near two stopped endpoints.
The latter calculation uses fixed normal sections when \(b/a\) is
near two. The remaining growing ranges have a strict curvature sign,
and a paired derivative bound supplies the upper complement.

Fix a preliminary \(0<\pi_{\rm pre}<\sqrt{12}\).
All comparison estimates are uniform for
\(\pi_{\rm pre}\le\pi\le\Pi\). The final
\(\pi_\flat\in(\pi_{\rm pre},\sqrt{12})\) will be chosen after
the comparison margin. For \(0<\pi\le\pi_\flat\),
Lemma~\ref{deg:lem-raw-jacobian} gives contraction after the final
coefficient shrink.

In \(x=gk\xi\) the exact polynomial and system are
\begin{equation}
 \begin{split}
 P(\xi)&=1-\pi\xi+3\xi^2+k\xi^3
       =(\mathfrak c+k\xi)\{(\xi-d)^2-\Delta\},\\
 \dot\xi&=-g\{P+Q_+Y\},\qquad
 \dot Y=-Y\{\xi B_++Ckg^3(1+k\xi)Y\},
 \end{split}
 \label{deg:quadratic-field}
\end{equation}
where
\[
 \mathfrak c=3+2kd,\qquad
 \mathfrak c(d^2-\Delta)=1,\qquad
 \pi=2\mathfrak cd-k/\mathfrak c,\qquad
 B_+=1+g(2+k\xi).
\]
Also \(Q_+=1-Ag(1+k\xi)+Cg^2(1+k\xi)^2\), and
the physical coordinates are \(x=gk\xi\), \(z=g^3k^2Y\).
At \(k=0\), \(d=\pi/6\) and
\(\Delta=(\pi^2-12)/36\). The curves \(Y=\ell_0P\) are transverse
by the analogue of \eqref{deg:compact-lip-speed}, and the reciprocal
raw equation \eqref{deg:compact-raw-corridor} joins them to the fixed
sections. On a provisional compact coefficient block about the fold,
and the real block up to \(\pi=\Pi\),
\(a,a^{-1},d,d^{-1},\mathfrak c,\mathfrak c^{-1}\) are bounded
where the corresponding real root is defined.
The original cocycles of Corollaries~\ref{deg:cor-positive-compact-suppliers}
and \ref{deg:cor-negative-inward-common-centre} give a finite
zeroth corrected-action bound \(K_0\). Fix \(r_c>0\) and a finite
action threshold greater than
\(K_0+\sup\widetilde A_-(r_c)+1\).
For a fixed small \(\delta_c>0\), the real factorization gives
\[
 P(v)\le C(a-v),\qquad vB_+(v)\ge c>0,
 \quad a-\delta_c\le v\le a-X_{\rm up}.
\]
Thus \(J_+(a-X_{\rm up})\ge c\log(\delta_c/X_{\rm up})\),
uniformly as \(b-a\to0\), and the fold has stronger
inverse-distance divergence. Choose \(X_{\rm up}=X>0\) to pay
this threshold. Choose the retained ghost width \(\Delta_0>0\)
sufficiently small afterwards so that the same margin holds on
\(-\Delta_0\le\Delta<0\). Then choose \(\pi_\flat\) with
\(0<12-\pi_\flat^2<18\Delta_0\).
Retain \(|\Delta|\le\Delta_0\) and the compact real block up to
\(\pi=\Pi\). Outside this band on the root-free side,
\(\pi<\pi_\flat\) after shrinking \(k\), so it is already a
contraction. In the remainder of the proof,
\begin{equation}
 \pi_\flat\le\pi\le\Pi.
 \label{deg:quadratic-pi-lower}
\end{equation}
Put \(u_{\rm th}=a-X\), \(a=d-\sqrt\Delta\), on the real/fold
side and \(u_{\rm th}=d-X\) on the ghost side.
The comparison \(u_*=a-X/2\), respectively \(d-X/2\), uses
this same \(X\) in Lemma~\ref{deg:lem-positive-quadratic-paired-variation}.
The independent degenerations are \(g\to0\), \(k\to0\), and the
positive root gap.

Fix a translated raw half-width \(0<d_q<\delta/2\), small enough for the
parabolic corridor estimates below, and shrink the coefficient neighborhood
so that \(|\gamma|<d_q/2\).  The cuts \(x=\pm d_q\), equivalently
\(v=\gamma\pm d_q\), lie on the corresponding sides of \(v=\gamma\) and
strictly within the strip bounded by the fixed sections
\(\Sigma_\pm=\{v=\pm\delta\}\).  The uniformly regular segments from these
cuts to \(\Sigma_\pm\), and their inverses, are included in the actual
section adapters used below.

For real or double roots the positive lip interval ends at \(a\).
For \(\Delta<0\), choose \(0<d_{\rm in}<d_0<d_q/2\) and restrict to
the component containing \(u_{\rm th}\) of the preimage of the
enlarged input and output section intervals, with
\[
 u\ge u_{\rm th},\qquad x=gku<d_0.
\]
Take smaller nested intervals with \(x<d_{\rm in}\).
Indeed \(z=\ell_0(x^3+3gx^2-\pi kg^2x+k^2g^3)\) on the lip;
for \(d_{\rm in}\le x\le d_q\) this is separated from the
organizing arc when \(\ell_0d_q\) is small.  The localization lemma
and the common transverse sign therefore place every actual low-height
cycle in the smaller interval.  On the positive continuation use the
relative outer cut \(x=L gku\), and on the negative continuation
\(x=-LgR\), with \(L>1\) fixed large.  At these cuts
\(\mathcal W=1/2-1/(2L^2)+o(1)\); the corridor equation carries
both traces to \(x=\pm d_q\), even across an axis root at positive
height.  This specifies the actual interval used below.  Its endpoint
on a root-free fibre is a finite section boundary, not \(u=\infty\).

Put \(R=kr\),
\[
 \widehat P(R)=k^2P(-R/k)
  =(\mathfrak c-R)\{(R+kd)^2-k^2\Delta\}.
\]
Here \(B_-(R)=1+g(2-R)\).  Define the two actions on fixed inner
cuts by \(J_+'=uB_+/P(u)\) and
\(A_-'=rB_- /P(-r)\), and put
\(\mathcal D_+=P(u)(uB_+)^{-1}\partial_u\).
In either negative coordinate,
\(\mathcal D_-= \widehat P(R)(RB_-)^{-1}\partial_R
=P(-r)(rB_-)^{-1}\partial_r\).
Retaining both logarithmic terms in the moving-lip matching gives
\[
 \widetilde J_+=J_+-g\log P,\qquad
 \widetilde A_-=A_--g\log P(-r),\qquad
 |\widetilde J_+-\widetilde A_-|\le K_0 .
\]
Here \(K_0\) is independent of \(g,k\), and
\(\widetilde A_-(r)\ge c\log(1+r)-C\).
The comparison action was chosen above
\(K_0+\sup\widetilde A_-(r_c)\).
Every retained low-height input then has \(r\ge r_c\), while the
comparison output itself is bounded above independently of \(k\).
We define \(\varrho\) by the exact relative variation
\[
 R_u=\frac{uB_+\widehat P(R)}{RB_-P(u)}(1+\varrho)>0 .
\]
On the growing negative tail let
\(\mathcal D_-=\widehat P(R)(RB_-)^{-1}\partial_R\).  Its logarithmic
action derivative is
\begin{equation}
 H_-={2\mathfrak c-3R\over1+g(2-R)}+O(r^{-1}+k),\qquad
 |\mathcal D_-^jH_-|\le C(R+r^{-1}),\quad j\ge1.
 \label{deg:quadratic-negative-action}
\end{equation}
In a fixed simple-root collar at \(R=\mathfrak c\), the rational
function \(H_-\) and its required derivatives are bounded, while
\(\mathcal D_-\) contains the factor \(\mathfrak c-R\).
Consequently \(\mathcal D_-^jH_-=O(\mathfrak c-R)\) for \(j\ge1\),
and \(dA_-\asymp dR/(\mathfrak c-R)\) there.  Together with the
growing-tail estimate this gives
\(\int|\mathcal D_-^jH_-|\,dA_-\le C\).  In the Bell densities
\eqref{deg:bell-density}, the factor \(P\) cancels with the action
change of variable.  The factor \(\mathfrak c-R\) is retained in the
exact logarithmic cocycle and the last segment is treated in the fixed
simple-root coordinate.  No power \(k^{-C/g}\) is generated.
More explicitly, on \(r_c\le r\le(\mathfrak c-\ell_-)/k\),
\(\mathcal D_-=V(kr,r^{-1})(kr\partial_{kr}
-r^{-1}\partial_{r^{-1}})\), with \(V\) a positive unit.
The integral of \(kr+r^{-1}\) against \(dA_-\asymp dr/r\)
is bounded on this interval.  The remaining fixed collar
\(\mathfrak c-2\ell_-<R<\mathfrak c\) has normal rate
\(RB_-\) bounded below and a simple factor \(\mathfrak c-R\).
The two estimates overlap at a fixed \(R\)-cut, where the
Leibniz boundary terms cancel.  No action integral crosses that root.

Only the root-free case can have \(u\to\infty\).  We first treat the
joint quadratic corner, where both radial action speeds are bounded.
Put \(q=u^{-1}\), \(S=ku\), and \(t=r^{-1}\).  Choose a small
\(\delta_q>0\), fix \(0<\delta<\delta_q/2\), and choose \(M\) large enough for
\(M\delta\gg\Pi\) and the action gap below.  All these constants
precede the choices of \(g,k\).  The homogeneous quadratic field has
first integral
\[
 Y^{-\lambda_g/2}\left(\kappa_gY-{\xi^2\over2}\right),
 \qquad \lambda_g={12g\over1+2g},\qquad
 \kappa_g=\frac{g(1-Ag+Cg^2)}{1-4g}.
\]
Its radial exponent is the same at the two fixed raw sections.  On
\(u,r\ge2M\), \(S,R\le\delta_q\), the exact ratio and variational
equations, after this power is extracted, give, with
\[
 h_+={P(u)\over u^2B_+(u)},\quad
 h_-={\widehat P(R)\over R^2B_-(R)},\quad
 \theta_+=u^{-1}+ku,\quad \theta_-=r^{-1}+kr,
\]
\begin{equation}
 \begin{split}
 \sum_{i+j\le2}|\mathcal D_+^i\mathcal D_-^j\varrho|
 &\le\varepsilon_g(\theta_++\theta_-),\\
 \mathcal D_\pm\log\phi_\pm'&=-\lambda_gh_\pm+G_\pm,\\
 |G_\pm|+|\mathcal D_\pm G_\pm|&\le\varepsilon_g\theta_\pm,
 \end{split}
 \qquad \varepsilon_g=Cg(1+|\log g|)^4.
 \label{deg:quadratic-weighted-estimates}
\end{equation}
The two \(-\lambda_gh_\pm\) terms come from the same homogeneous
power and must be combined before estimating the curvature.
These bounds hold in the joint corner; the range \(S\to\infty\),
where \(h_+\asymp1+S\), will instead use radial derivatives.
Lemma~\ref{deg:lem-joint-inward-product-direct} gives the first
line on the product of the two separate actual domains, including
the mixed second row; restriction to their original central-image
intersection gives the matched graph. Lemma~\ref{deg:lem-quadratic-outer-jets}
gives the section rows with the stronger factor \(Cg(1+|\log g|)\),
on each admitted raw image.
The negative observer includes the divided closing of
Lemma~\ref{lem:closing}; its ordinary \(C^3\) logarithmic-slope
bounds apply on those original buffered images, with its small gain
and \(m\) fixed under the phase derivative.

We shall also use absolute negative-side estimates on the larger range
\(r\ge r_c\), \(0<R<\mathfrak c\).  Write the exact logarithmic
matching as
\[
 J_+(u)-A_-(r)
 =g\{\log P(u)-\log P(-r)+\mathcal R_+(u)-\mathcal R_-(r)\}.
\]
Coefficient-only constants may be included in either \(\mathcal R_\pm\).
The two functions are obtained by continuing each actual lip to the same
central section.  Thus their separated extension and its derivatives
refer to the actual orbit.  The complete negative inward action remainder and the negative
outward embedding give
\begin{equation}
 |\mathcal D_-^j\mathcal R_-|\le Cg\quad(1\le j\le3),\qquad
 |\mathcal D_-^j\log\phi_-'|\le C\varepsilon_g\quad(j=1,2).
 \label{deg:quadratic-negative-absolute}
\end{equation}
Here and below primes on \(\phi_-\) are in the label \(-R^2/2\).
Changing from \(-r^2/2\) only changes its frozen affine slope.
The first estimate is Corollary~\ref{deg:cor-negative-inward-common-centre}
at the original centre, retaining the initial polynomial and every
coefficient-only action constant. In particular,
\[
 N_-=-g\mathcal D_-\log Z_-
 =1-gH_--g\mathcal D_-\mathcal R_-,\qquad
 |\mathcal D_-^j(N_--1+gH_-)|\le Cg^2,\quad0\le j\le2.
\]
Since \(H_-\) is bounded, \(N_-\ge1/2\). The second estimate is
Corollary~\ref{deg:cor-negative-absolute-outer}, with the stronger
\(Cg(1+|\log g|)\) bound, including original-buffer admission,
\(p\)-weighted initial variations, both moving endpoints and the
ordinary \(C^3\) divided closing. Neither estimate enlarges the
intersection of actual central images.

For \(\mathcal T(-u^2/2)=-R^2/2\) and
\(\Pi=\phi_-\circ\mathcal T\circ\phi_+^{-1}\), put
\(B_{\rm rat}=B_+/B_-\) and
\(\mathcal V=\mathcal D_++(1+\varrho)\mathcal D_-\).
Lemma~\ref{neg:lem-completed-curvature}, with
\(v=R\), \(P_-=\widehat P\), and \(W=R+S\), gives
\eqref{deg:weighted-exact-curvature-chain} for the weighted curvature
\begin{equation}
 \mathscr K=\phi_+'\frac{P(u)}{R+S}\frac{\Pi''}{\Pi'}.
 \label{deg:quadratic-full-chain}
\end{equation}
The derivative \(d/du\) follows the matched graph.  In the joint
corner, exact logarithmic differentiation and the preceding estimates
give
\begin{equation}
 \phi_+'{P(u)\over R+ku}{\Pi''\over\Pi'}
 =F_\lambda(u,R)+E(u,R),
 \label{deg:quadratic-curvature}
\end{equation}
where, with \(q=u^{-1}\) and \(t=r^{-1}\),
\begin{equation}
 F_\lambda=(3-\lambda_g)-{N_\lambda\over uR},\qquad
 N_\lambda=(1-\lambda_g)\pi-\lambda_g(t-q),
 \label{deg:quadratic-principal}
\end{equation}
and, in the joint corner \(u,r\ge2M\), \(S,R\le\delta_q\),
\begin{equation}
 |E|+|(\mathcal D_++(1+\varrho)\mathcal D_-)E|
 \le C\varepsilon_g\left(1+{1\over uR}\right).
 \label{deg:quadratic-error}
\end{equation}
To verify \eqref{deg:quadratic-principal}, use
\[
 {P'(u)/u-\widehat P'(R)/R\over R+ku}=3-{\pi\over uR}
\]
and
\[
 {\widehat P/R^2-P/u^2\over R+ku}
 ={ku-R\over u^2R^2}+{\pi\over uR}-1.
\]
The matching rows of Lemma~\ref{deg:lem-joint-inward-product-direct}
apply on the product domain to
\(N_\pm=1-gH_\pm-g\mathcal D_\pm\mathcal R_\pm\),
\(n=1+\varrho=N_+/N_-\), with \(N_-\ge1/2\).

To estimate the curvature error and its derivative along the match,
put \(W=R+S\), \(\mathcal V=\mathcal D_++n\mathcal D_-\), and
\(A_\pm=B_\pm(H_\pm-\lambda_gh_\pm)\). Direct polynomial calculation gives
\begin{align*}
 A_+&=(6-3\lambda_g)+(\lambda_g-1)\pi q
                 +(3-\lambda_g)S-\lambda_gq^2,\\
 A_-&=(6-3\lambda_g)+(1-\lambda_g)\pi t
                 -(3-\lambda_g)R-\lambda_gt^2.
\end{align*}
Since \(B_+-B_-=gW\), the exact error decomposition is
\begin{align*}
 E_B&=-gA_-/B_--(B_+/W)\mathcal V\log(B_+/B_-),\\
 E_n&=-(B_+/W)\{\varrho(H_--\lambda_gh_-)+\mathcal V\log n\},\\
 E_s&=(B_+/W)(G_+-nG_-),\qquad E=E_B+E_n+E_s.
\end{align*}
The first two action rows of \(h_\pm,B_\pm,H_\pm\) are bounded.
Thus \(|E_B|+|\mathcal VE_B|\le Cg\). The total-two matching row gives
the required bound on \(E_n\), including \(\mathcal V^2\log n\).
For the section term, retain the exact derivative
\[
 \mathcal V(G_+-nG_-)
 =\mathcal D_+G_+-(\mathcal Vn)G_--n^2\mathcal D_-G_-.
\]
Lemma~\ref{deg:lem-quadratic-outer-jets} bounds these terms with the
stronger factor \(Cg(1+|\log g|)\). Finally
\((q+S+t+R)/W=1+1/(uR)\), and \(|\mathcal V\log W|\le C\).
These identities give \eqref{deg:quadratic-error} on the actual
matched joint graph. The other coefficient ranges are treated below.

On the joint corner \(q\) and \(t\) are small.  The lower bound
\eqref{deg:quadratic-pi-lower} and a further common shrink give
\[
 N_\lambda=(1-\lambda_g)\pi-\lambda_g(t-q)
 \geq\frac14\pi_\flat .
\]
Put \(z=uR\) and \(a_0=3-\lambda_g\). At a zero of the actual
\(\mathscr K=F_\lambda+E\), write
\(|E|\le\delta_g(1+1/z)\), with \(\delta_g\to0\).
Then
\[
 \frac{N_\lambda-\delta_g}{a_0+\delta_g}
 \le z\le
 \frac{N_\lambda+\delta_g}{a_0-\delta_g}.
\]
All actual curvature zeros therefore lie in a fixed positive compact
\(z\)-interval. The exact derivative there is
\[
 \mathcal VF_\lambda
 =\frac{N_\lambda(h_++nh_-)
          +\lambda_gh_+q-\lambda_gnh_-t}{z}\ge c>0,
\]
after \(M\) is large and \(g\) small. The differentiated error is
absorbed on this compact interval. If
\(w=\phi_+(-u^2/2)\) is the actual affine input, then
\(\mathcal Vw=-P(u)\phi_+'/B_+<0\). Because \(\mathscr K\) is a positive
multiple of \(\Pi''/\Pi'\), every curvature zero on this matched joint
graph has \(\Pi'''(w)<0\). To obtain this orientation on the
entire low-height interval, we next treat the mixed and remaining
growing ranges, followed by the bounded positive range.
There is also a mixed range \(u\ge2M\), \(r_c\le r\le2M\).
If \(S\ge\delta\), the action comparison below already forces
\(R\ge R_0(\delta)>0\), which contradicts \(R\le2Mk\).
On \(S\le\delta\), the positive half of the joint-corner calculation
still applies: its action speed is bounded.  The negative end is now
compact in \(r\), and \eqref{deg:quadratic-negative-absolute} supplies
its absolute estimates.  In \eqref{deg:weighted-exact-curvature-chain} these give
\(\mathscr K=3-\pi/(uR)+E_M\), with
\(|E_M|\le C_M\varepsilon_g/(R+S)\).  Relative to
\(1/(uR)\), it is bounded by
\(C_M\varepsilon_g\,ur/(u+r)\le2MC_M\varepsilon_g\).
A zero would require \(uR\ge c_M>0\), and hence
\(ku=uR/r\ge c_M/(2M)\).  The action comparison below applies
already for \(r\ge r_c\); with this last fixed threshold in place
of \(\delta\), it forces \(R\ge R_0>0\), contrary to
\(R=kr\le2Mk\).  Thus the mixed range contains no zero.
The remaining corners introduce none.  If \(u\ge2M\),
\(R\ge\delta\), and \(S\le\delta\), the same positive estimates
and \eqref{deg:quadratic-negative-absolute}, including its simple-root
collar, give an error \(O_\delta(\varepsilon_g)\) in
\eqref{deg:weighted-exact-curvature-chain}.  The exact polynomial principal part
is \(3-\pi/(uR)\); the choice \(M\delta\gg\Pi\) leaves a
positive margin.  For the last range \(S\ge\delta\), fix
\(0<\overline R<\tfrac12\inf\mathfrak c\).
For \(r\ge r_c\) and \(R\le\overline R\), subtraction of the exact
logarithmic cocycles gives
\begin{equation}
 \widetilde J_+(u)\ge {1-4g\over3}\log{1\over k}-C_\delta,
 \qquad
 \widetilde A_-(r)\le {1-4g\over3}
   \left(\log{1\over k}+\log R\right)+C_\delta.
 \label{deg:quadratic-large-S-action}
\end{equation}
The upper bound for \(\widetilde A_-\) is used only on this interval,
which stays away from its logarithmic singularity at
\(R=\mathfrak c\).  Write \(c_g=(1-4g)/3\), and shrink \(g\) so
that \(c_g\ge c_*>0\).  Since
\(|\widetilde J_+-\widetilde A_-|\le K_0\), the displayed inequalities
give \(c_g\log R\ge-K_0-2C_\delta\).  Together with the alternative
\(R>\overline R\), this proves the uniform bound
\[
 R\ge R_0(\delta):=
 \min\{\overline R,\exp[-(K_0+2C_\delta)/c_*]\}>0.
\]
Consequently \(uR\ge\delta R_0(\delta)/k\to\infty\).
We prove the required estimate on this range in the radial derivative
\(E=u\partial_u\), keeping the conversion to action derivatives
explicit.  Set
\[
 a=gS=gku,\qquad L_a=1+|\log a|,\qquad
 h_+=\frac{3-\pi/u+u^{-2}+S}{1+2g+gS}.
\]
For \(S\ge\delta\), \(a<d_0\), and \(k\) small,
\begin{equation}
 h_+\asymp_\delta1+S,\qquad
 |E^jh_+|\le C_\delta h_+\quad(j=1,2),\qquad
 \frac{h_+}{R+S}\le C_\delta,
 \qquad \mathcal D_+=h_+E.
 \label{deg:quadratic-large-S-speed}
\end{equation}
In particular the first bound in
\eqref{deg:quadratic-weighted-estimates} is not used here.

Corollary~\ref{deg:cor-positive-large-S-inward}, with
\(k_0\le\delta/(2M)\), gives the actual same-centre Euler-three rows
\begin{equation}\label{deg:quadratic-large-S-action-jets}
 |E^j\mathcal R_+|\le C_\delta a,\qquad1\le j\le3.
\end{equation}
Its strict ghost domain retains the moving initial polynomial,
terminal terms and full compact multiplier; the real/fold through
interval is empty here. The action conversion remains \(\mathcal D_+=h_+E\).

Differentiating the separated action matching gives exactly
\[
 1+\varrho=
 \frac{1-gH_+-g\mathcal D_+\mathcal R_+}
      {1-gH_--g\mathcal D_-\mathcal R_-},\qquad
 H_+=\frac{P'(u)}{uB_+(u)}.
\]
Both denominators remain positive.  Since
\(g|E^jH_+|\le C_\delta a\), \(gh_+\le C_\delta a\),
\eqref{deg:quadratic-negative-absolute} and
\eqref{deg:quadratic-large-S-action-jets} imply
\begin{equation}
 |\varrho|+|E\varrho|\le C_\delta a,\qquad
 |\mathcal D_-\varrho|\le C_\delta g,\qquad
 |\mathcal V\log(1+\varrho)|
       \le C_\delta(h_+a+g).
 \label{deg:quadratic-large-S-relative}
\end{equation}
Thus an action derivative of \(\varrho\) need not be small before
division by \(R+S\).

Choose \(d_0<a_0\) from
Corollary~\ref{deg:cor-positive-large-S-outer}. Its actual positive
raw graph, with \(k_0\le\delta/(2M)\), original observer and frozen
\((gk)^2\) first gain, gives the stronger \(C_\delta aL_a\) bound;
hence
\begin{equation}\label{deg:quadratic-large-S-section-jets}
 |E^j\log\phi_+'|\le C_\delta aL_a^4,\qquad j=1,2.
\end{equation}

We now substitute in \eqref{deg:weighted-exact-curvature-chain}.  Its polynomial
term has principal part \(3-\pi/(uR)\).  Since
\(B_{\rm rat}-1=g(S+R)/B_-\) and
\(\widehat P_R/R\) is bounded for \(R\ge R_0\), the remaining
part of that term is \(O_\delta(a)\).  The second term is exactly
\[
 -\frac{gS h_+}{R+S}
 -\frac{gB_+\widehat P(1+\varrho)}{RB_-^2(R+S)},
\]
and is again \(O_\delta(a)\).  The third and the positive section
term are bounded by \eqref{deg:quadratic-large-S-speed},
\eqref{deg:quadratic-large-S-relative}, and
\eqref{deg:quadratic-large-S-section-jets}.  For the complete negative
section, including its root collar and finite closing, use
\eqref{deg:quadratic-negative-absolute} and
\[
 \frac{g(1+|\log g|)^4}{R+S}\le C_\delta aL_a^4.
\]
Indeed \(a=gS\), \(S\ge\delta\), and
\((1+\log S)^4/S^2\) is bounded for \(S\ge1\); the interval
\(\delta\le S\le1\) is compact.  We have proved
\begin{equation}
 \mathscr K=3-\frac{\pi}{uR}+E_S,\qquad
 |E_S|\le C_\delta a(1+|\log a|)^4.
 \label{deg:quadratic-large-S-curvature}
\end{equation}
First decrease \(d_0\) so that this error is small on \(a<d_0\),
then decrease \(g,k\).  The lower bound on \(uR\) gives strict
positivity on the entire large-\(S\) range.

It remains to justify the bounded positive range and its overlap with
the preceding estimates.  On \(u\le3M\), Corollary~\ref{deg:cor-positive-compact-suppliers},
prepared on \(u<4M\), supplies
\( |\mathcal D_+^j\mathcal R_+|\le C_Mg\), \(1\le j\le3\), and
\( |\mathcal D_+^j\log\phi_+'|\le C_Mg(1+|\log g|)\), \(j=1,2\).
Thus \(N_+=1-gH_+-g\mathcal D_+\mathcal R_+\ge1/2\), with
\(|\mathcal D_+^j(N_+-1+gH_+)|\le C_Mg^2\), \(j\le2\),
on its original real/fold interval or strict ghost preimage.

On bounded moving-lip abscissae the ratio
equation, the growing negative estimate
\eqref{deg:quadratic-negative-action}, and the full section maps give
\begin{equation}
 \mathscr K=3-\frac{\pi}{uR}+E_{\rm lip},\qquad
 |E_{\rm lip}|+|\mathcal D_+E_{\rm lip}|
                  +|\mathcal D_-E_{\rm lip}|
       \le\frac{C\varepsilon_g}{R+k},
 \quad
 \mathscr K=\phi_+'\frac{P(u)}{R+ku}\frac{\Pi''}{\Pi'}.
 \label{deg:quadratic-bounded-lip-curvature}
\end{equation}
These action rows consume the bounded positive and absolute
negative suppliers on their separate actual images, with both
moving initial factors, endpoint rows and physical gains retained.

On a real-root fibre factor
\[
 P(\xi)=(\mathfrak c+k\xi)(\xi-a)(\xi-b),\qquad
 0<a\le b,\quad
 \mathfrak c=3+k(a+b),\quad
 \mathfrak cab=1,\quad
 \pi=\mathfrak c(a+b)-kab.
\]
Here \(a,a^{-1},b,\mathfrak c,\mathfrak c^{-1}\) are bounded.
At the simultaneous stopped endpoint \((u,R)=(a,\mathfrak c)\),
the limiting principal value at \(g=k=0\) is
\[
 3-\frac{\pi}{a\mathfrak c}=2-\frac ba.
\]
The endpoint value vanishes at \(b/a=2\); this is distinct from
the double-root condition \(\pi=\sqrt{12}\) at \(k=0\).
Outside a fixed neighborhood of \(b/a=2\), this value has a
nonzero margin.  A possible curvature zero then has
\(R\ge c_F>0\) and \(P(u)+\widehat P(R)\ge c_0>0\).
Since
\[
 \frac{d}{du}
 =\frac{uB_+}{P(u)}
       \{\mathcal D_++(1+\varrho)\mathcal D_-\},
 \qquad
 \frac{uB_+}{P(u)}\le C(1+R_u)
\]
on this zero neighborhood, the error derivative is at most
\(C\varepsilon_g(1+R_u)\), whereas
\(d(3-\pi/(uR))/du\ge c(1+R_u)\).
For \(R\to0\), the principal term instead dominates the error
relatively.  This proves the crossing assertion except in a fixed
joint-root neighborhood of \(b/a=2\).  The double root \(a=b\)
belongs to the separated case.  The adjacent root-free block has the
same endpoint margin: the arctangent action gives its continuation,
and no passage beyond \(a=d\) is added on \(\Delta=0\).

Near \(b/a=2\), the endpoint principal value can vanish while both
action speeds tend to zero. We therefore need errors whose ordinary
derivatives are small in the two independent root distances.
The root gap \(b-a\) is a positive unit in this range.
Use the actual fixed normal sections constructed in
Lemma~\ref{deg:lem-joint-root-physical-resection}, in the positive
\((\xi,gY)\) and negative \((R,gk^2Y)\) root boxes.  Their strong footpoints
are \(u=a-h\) and \(R=\mathfrak c-t\), with \(h,t>0\).
The retained interval is the intersection of the actual matched graph
with \(h<h_0\), \(t<t_0\), where \(h_0,t_0\) are fixed first.
It is one terminal interval because \(t'(h)>0\).
The normal rates in these two boxes are positive units, independently
of \(g,k\).  In this fixed-root calculation, regard \(A_-\) as
a function of \(R\); the identity
\(rB_-\,dr/P(-r)=RB_-\,dR/\widehat P(R)\) preserves its value.
Lemma~\ref{deg:lem-joint-root-physical-resection} supplies the
pole subtractions with residues \(\mathfrak a_\pm\) and uniform
\(C^3\) regular parts. All coefficient-only action constants,
including possible \(\log k\), remain fixed. The normal-footpoint
quotient is \(\nu_{\rm fp}=\mathfrak a_+/\mathfrak a_-\), without
an added \(g\) in either residue.
The exact residue identity is
\begin{equation}
 3-{\pi\over a\mathfrak c}
 -{g\widehat P_R(\mathfrak c)\over
        \mathfrak cB_-(\mathfrak c)}
 ={b-a\over a}\left({\mathfrak a_+\over\mathfrak a_-}-1\right).
 \label{deg:quadratic-residue-identity}
\end{equation}
This identifies the endpoint value of the actual weighted curvature;
it does not assert that every nearby curvature zero lies on the
coefficient resonance itself.

On the actual matched intersection, the same lemma gives
\eqref{deg:joint-root-normal-weighted-matching},
\eqref{deg:joint-root-normal-matching-rows}, and
\eqref{deg:joint-root-normal-outer-rows}. Its coefficient memory
\(C_{g,k}\) and both homogeneous gains are retained without size
or reciprocal bounds; the separated matching errors vanish with
\(h,t\), and \(|\partial_h\varrho|+|\partial_t\varrho|\le C\varepsilon_g\).

Put \(E_*=-g\widehat P_R(\mathfrak c)/
(\mathfrak cB_-(\mathfrak c))\).
The axis correction in the exact curvature identity is
\[
 E^{\rm ax}=
 -\frac{gkP}{u(R+ku)B_+}
 -\frac{g\widehat P_R}{RB_-}
 -\frac{gB_+\widehat P}{R(R+ku)B_-^2}.
\]
It has endpoint value \(E_*\).
The remaining matching and section terms, before estimating, are
\[
 \begin{split}
 E^\varrho={}&
 -\frac{B_+\varrho(\widehat P_R-\widehat P B_{-,R}/B_-)}
                    {RB_-(R+ku)}
 -\frac{B_+\{\mathcal D_+\varrho+
                    (1+\varrho)\mathcal D_-\varrho\}}
                    {(R+ku)(1+\varrho)},\\
 E^{\rm sec}={}&
 \frac{B_+}{R+ku}
  \{\mathcal D_+\log\phi_+'-
                 (1+\varrho)\mathcal D_-\log\phi_-'\}.
 \end{split}
\]
Write \(E_0=E^{\rm ax}-E_*+E^\varrho+E^{\rm sec}\).
Here \(P\asymp h\), \(\widehat P\asymp t\), and all displayed
denominators are positive units with bounded ordinary endpoint derivatives.
The separated quotient calculation in
Lemma~\ref{deg:lem-joint-root-physical-resection} bounds the ordinary
\(h,t\) derivatives of \(\varrho\), \(\mathcal D_+\varrho\)
and \(\mathcal D_-\varrho\) by \(C\varepsilon_g\).
For the section terms, division of the second rows in
\eqref{deg:joint-root-normal-outer-rows} by
\(h\gamma_+\), respectively \(t\gamma_-\), gives
\[
 |\partial_h(\mathcal D_+\log\phi_+')|
 +|\partial_t(\mathcal D_-\log\phi_-')|
 \le C\varepsilon_g .
\]
The axis correction satisfies
\(E^{\rm ax}-E_*=O(g(h+t))\), with ordinary endpoint derivatives
\(O(g)\). Hence the product rule gives
\( |\partial_hE_0|+|\partial_tE_0|\le C\varepsilon_g\)
on the separated product. Restricting to the actual \(t=t(h)\) yields
\begin{equation}
 \mathscr K=3-\frac{\pi}{uR}+E_*+E_0,\qquad
 |E_0|\le C\varepsilon_g(h+t),\qquad
 \left|\frac{dE_0}{dh}\right|
       \le C\varepsilon_g(1+t'(h)),
 \label{deg:quadratic-normal-curvature}
\end{equation}
Since
\[
 \frac{d}{dh}\left(3-\frac{\pi}{uR}\right)
 =-\frac{\pi}{u^2R}
  -\frac{\pi t'(h)}{uR^2}
 \le-c(1+t'(h)),
\]
the actual curvature is strictly decreasing at every possible zero,
even when both action rates tend to zero independently.

To join this interval to the moving lip, use the actual nonlinear
resection of Lemma~\ref{deg:lem-joint-root-physical-resection}.
Equation~\eqref{deg:joint-root-exact-lip-resection} retains the added
\(g\)'s from the two initial polynomials. Its positive-unit residues
give uniform inclusions between suitably nested fixed endpoint
neighborhoods; they do not give a linear comparison of the internal
root distances. Equation~\eqref{deg:joint-root-physical-composition}
identifies the same completed physical map after the two exterior
embeddings. Only positive affine changes of its physical input and
output remain, so its curvature zeros agree. Both inputs increase
with the corresponding positive endpoint distance. Hence every zero
has the same downward orientation on the overlap.
The bounded construction is used on \(u\le3M\).  On \(u\ge2M\),
the mixed range covers \(r\le2M\); for \(r\ge2M\), the remaining
alternatives are \(S,R\le\delta\), \(R\ge\delta,S\le\delta\),
and \(S\ge\delta\).  They are covered respectively by the joint
corner, the absolute negative-tail estimate, and
\eqref{deg:quadratic-large-S-curvature}.  The fixed thresholds are used
with the enlarged overlaps specified above, so these constructions
exhaust the prescribed interval. Every curvature zero has this same
strict downward orientation in the common physical input.

Put \(u_*=a-X_{\rm up}/2\) on the real/fold side and
\(u_*=d-X_{\rm up}/2\) on the ghost side. In the same original
physical input, Lemma~\ref{deg:lem-positive-quadratic-paired-variation}
gives \(0<\Pi_\theta'\le C_*\) on the entire admitted interval
\([\eta_{\rm comp},\eta_{\max}]\), where
\(\eta_{\rm comp}=gZ_+(u_*)\). The bounded comparison cocycle
gives
\(\eta_{\rm comp}<gZ_+(u_{\rm th})<\eta_{\max}\) after a common
coefficient shrink. Hence this interval overlaps the lower lip
interval strictly and retains every intervening input in the
original buffers. The common choice \(mC_*<1/2\) gives
\(F'>1/2\) on its bounded-height complement, uniformly through
\(\Delta=0\).
On \(\Delta=0\), the formal branch
\(u>d\) approaches the equilibrium in infinite time and is not included.
\end{proof}

\subsection{The admitted raw tail on a bounded positive quadratic fibre}

The growing curvature calculation uses \(gku<d_0\). On a ghost
fibre, the original physical interval can also contain first lips beyond
this divider. We keep that tail and compare its complete input and output
gains through their common central height. The negative gain is controlled
by the sign of its density during the long root dwell.

\begin{proposition}[The original bounded-positive ghost raw tail]
\label{bp-tail:prop-actual-raw-tail}
Fix the passive compact set, original raw sections, nested windows,
ordinary \(C^3\) observers and finite closing of Lemma~\ref{lem:closing}.
Fix \(\pi\in[\pi_\flat,\Pi]\) and the retained strict ghost band
\(\Delta(k,\pi)<0\) of \eqref{deg:quadratic-field}. The parameters
\(g,k>0\) approach zero independently. Keep the original master
\[
 F(x)=x-m\Pi_{\theta,m}(x),\qquad x=h-a_\theta,
 \qquad \Pi_{\theta,m}=\mathcal C_{\theta,m}\circ T_\theta(a_\theta+\cdot),
\]
and one positive unscaled cap
\(0<h_c<\min(H_{\rm cap},z_*)\), chosen before the final coefficient
and multiplier bounds. Its original carrier is
\begin{equation}\label{bp-tail:original-carrier}
 I_\theta=\{x:a_\theta+x\in J_\theta^{\rm out},\quad
                         0<H(x)=z_\theta(g;a_\theta+x)<h_c\}.
\end{equation}
Choose the fixed auxiliary \(d_0<a_0\) and its buffer \(a_b=d_0/2\)
inside the original chart, where \(a_0\) is the positive outward
reserve of Corollary~\ref{deg:cor-positive-large-S-outer}.
There are fixed \(g_{\rm tail},k_{\rm tail}>0\);
the full-carrier coefficient rectangle uses the minima of these bounds
and the native lower, paired and high bounds. Every actual lower input whose positive first
lip has raw abscissa \(a=gku\ge a_b\) satisfies
\begin{equation}\label{bp-tail:consumed-derivative}
                   0<\Pi_{\theta,m}'(x)\le C(d_0)g^3.
\end{equation}
This includes every original admitted \(a>d_0\), through all original
window, inverse-image, speed, stop and height truncations, and all
intervening nonroot inputs. The estimate holds uniformly in the independent
coefficient rectangle and \(0\le m\le m_0\). With the actual finite
channel \(K_{\rm fin}>0\), choose \(m_*\) last so the whole retained
raw tail has \(F'>1/2\). No new receiver or first-lip cutoff is imposed.
\end{proposition}

\begin{proof}
Put \(\epsilon=gk\), \(\eta=gY\), so
\(v-g=\epsilon\xi\), \(z=\epsilon^2\eta\). Use \eqref{deg:quadratic-field} in the original
\(\eta=gY\) coordinate, with \(P=1-\pi\xi+3\xi^2+k\xi^3\),
\(B=1+g(2+k\xi)\), \(Q=\ell(g+\epsilon\xi)\),
\(\ell(v)=1-Av+Cv^2\), and \(H_1=Ckg^2(1+k\xi)\).
On the ghost fibre \(P>0\) on the entire positive segment. In the
translated raw positive abscissa \(t=v-g\), define
\[
 \begin{gathered}
 M(t)=t^3+3gt^2-\pi kg^2t+k^2g^3,\quad B_0(t)=1+2g+t,\\
 Q_0(t)=\ell(g+t),\quad\zeta_0(t)=C(g+t),\quad D_0=M+Q_0z.
 \end{gathered}
\]
Thus \(M(\epsilon\xi)=\epsilon^2gP(\xi)\).
Let \(t_*=\delta-g\) be the actual original positive raw port.
The inverse of its incoming passage obeys
\begin{equation}\label{bp-tail:positive-raw-graph}
 z_t=f_0(t,z)=\frac{z[tB_0+\zeta_0z]}{D_0},\qquad
                         z(a)=\ell_0M(a).
\end{equation}
The lip height \(\ell_0\) is fixed using the original input-window
reserve. On \(a_b\le a\le t_*\), independently small \(g,k\) give
\(M(t)\asymp t^3\) on \([a,t_*]\). With \(Q_0\) a positive unit,
\(f_0\le Ct+Cz\); Gronwall gives \(z\le C\delta^2\) inside the
fixed upper chart boundary. The numerator is positive, so
\(z\ge\ell_0M(a)>0\). Hence this is an actual regular
outward flight. Its original input observer \(\chi_+\) is used only
on admitted port images. A nonadmitted zero-flight endpoint at \(a=t_*\)
is an auxiliary graph value and is not assigned an observer value.

Equation~\eqref{deg:scalar-polynomial-lip-slope}, with
\(\sigma=\epsilon\), \(K_l=K\), \(I_+=I\), and \(x_T=t_*\), gives
\begin{equation}\label{bp-tail:positive-full-gain}
 \Gamma_{+,a}'=-a\chi_+'(z_{\rm port})\frac{z_{\rm port}}{D_0(t_*,z_{\rm port})}
                     B_0(a)K_l(a)e^{I_+(a)}.
\end{equation}
The original transverse lip margin makes \(K_l\) a positive unit.
All admitted remaining factors have bounds and reciprocal bounds
depending on \(d_0\). For the original label \(b_u=-u^2/2\),
\begin{equation}\label{bp-tail:positive-raw-slope}
 -\Gamma_{+,a}'\ge c(d_0)a,\qquad
 c(d_0)\epsilon^2\le\Gamma_{+,b_u}'\le C(d_0)\epsilon^2.
\end{equation}
The exact input gain \(\epsilon^2=(gk)^2\) is retained.

We next obtain an absolute first-variation estimate for the positive
inward passage on the whole raw tail. Let
\(\eta_0^+(u)>0\) be its height at the original center \(\xi=0\).
For \(f=\eta(\xi B+H_1\eta)/(gP+Q\eta)\), put
\(W=\eta/(gP)\), \(H_+=P'/(\xi B)\),
\(\varpi=H_1gP/(\xi B)\). Moving-initial differentiation is exact:
\begin{equation}\label{bp-tail:positive-inward-variation}
 \begin{gathered}
 \partial_u\log\eta_0^+=-\frac{uB_+}{gP_+}\Lambda_+(u),\qquad
 \Lambda_+=\frac{K_l(a)}{1+\ell_0Q_+}
                  \exp\int_u^0\left(f_\eta-\frac f\eta\right)d\xi,\\
 f_\eta-f/\eta=\frac{\xi B}{gP}
                 \frac{W(\varpi-Q)}{(1+QW)^2}.
 \end{gathered}
\end{equation}
Indeed \(g\ell_0P'(u)-f(u,g\ell_0P(u))
=-\ell_0uB_+K_l(a)/(1+\ell_0Q_+)\); division by the actual
height cocycle produces the displayed logarithmic variation.
Fix \(v_i>0\) so that \(P\ge p_i>0\) on \([0,2v_i]\).
For \(\xi\ge v_i\) throughout the original raw chart,
\[
 gH_+=\frac{g(6-\pi/\xi+3k\xi)}B\le C(g+t),\qquad
 |\varpi|\le C_{v_i}\{(g+t)^3+g^3k\}.
\]
The fixed chart and lip choices therefore give \(Q-\varpi>0\)
and the same transverse margin at every running point.
The positive action clock \(dq=-\xi B\,d\xi/(gP)\) gives
\[
 W_q=W\left[-\frac{1+\varpi W}{1+QW}+gH_+\right],\qquad W(0)=\ell_0.
\]
The braces increase with \(W\); at \(\ell_0\) they equal
\(-K_l/(1+\ell_0Q)\le-c_1<0\). Thus
\(W\le\ell_0e^{-c_1q}\) down to \(\xi=v_i\).
In raw coordinates this also gives
\(z=M(t)W\le\ell_0M(t)\le C\ell_0\delta^3\), keeping the
entire noncentral inward arc in the original chart.
The noncentral integral in \eqref{bp-tail:positive-inward-variation}
is exactly \(\int W(Q-\varpi)/(1+QW)^2\,dq\), and is bounded.

Every \(a\ge a_b\) has \(u>2v_i\) after independent coefficient
shrinks. The fixed segment \([v_i,2v_i]\) gives
\(A_{\rm in}=\inf\int_{v_i}^{2v_i}\xi B/P\,d\xi>0\).
Consequently its entering height satisfies
\(\eta_i\le Cg e^{-c_1A_{\rm in}/g}\).
On the remaining slab \([0,v_i]\), the linear coefficient
\(\xi B/(gP)\) is nonnegative, and the graph equation has
\(f\ge-Ckg\,\eta^2\). In the decreasing \(\xi\)-direction a
first-exit comparison therefore keeps \(0<\eta\le2\eta_i\).
The denominator remains at least \(gp_i/2\), so the actual
arc reaches \(\xi=0\) in finite positive time. The regular formula
\[
 f_\eta-f/\eta
 =\frac{\eta(H_1gP-\xi BQ)}{(gP+Q\eta)^2}
\]
bounds the remaining logarithmic variational integral by
\(C\eta_i/g^2\), which is exponentially small. This uses only the
original central passage, with no additional observer or endpoint.
Consequently
\begin{equation}\label{bp-tail:positive-inward-absolute}
 0<c\le\Lambda_+,\Lambda_+^{-1}\le C,
 \qquad -\partial_a\log\eta_0^+
             =\frac{aB_0(a)}{M(a)}\Lambda_+>0.
\end{equation}
This proves monotonicity and absolute first variation for every actual
raw-tail input, with no restriction \(a<d_0\).

The negative inward supplier is already on its entire strict interval.
Write \(R=kr\),
\[
 P_-(r)=1+\pi r+3r^2-kr^3,\qquad
 \widehat P(R)=k^2P_-(R/k),\qquad B_-=1+g(2-R).
\]
Corollary~\ref{deg:cor-negative-inward-common-centre} constructs the
actual original-center height \(\eta_0^-(r)=gZ_-(r)\) and retains
\(\log(\ell_0P_-)-A_-/g+\mathcal R_-\). Its exact derivative is
\begin{equation}\label{bp-tail:negative-inward-absolute}
 \begin{gathered}
 -\partial_r\log\eta_0^-=\frac{rB_-}{gP_-}\Lambda_-,\qquad
 \Lambda_-=1-gH_--g\mathcal D_-\mathcal R_-,\qquad
                         1/2\le\Lambda_-\le3/2,\\
 \mathcal D_-=\frac{P_-}{rB_-}\partial_r,\qquad
 r\ge r_c,\quad 0<R<\mathfrak c.
 \end{gathered}
\end{equation}
This includes the arbitrarily close negative root, its actual endpoint
power and all central continuation factors. In particular
\(\eta_0^-\downarrow0\) at \(R\uparrow\mathfrak c\).

Matching does not assume an image intersection or an action estimate
outside its original domain. The original paired comparison has an
actual negative match \(r_*\ge r_c\) at \(u_*\).
Put \(u_b=a_b/(gk)\). The native separate inward estimates give
\(\eta_0^+(u_b)<\eta_0^+(u_*)=\eta_0^-(r_*)\), since their entire
segment has \(a\le a_b<d_0\). Negative
\eqref{bp-tail:negative-inward-absolute} and its zero endpoint
therefore give a unique \(r_b\ge r_*\). This is a match of the
separate actual centre maps, before any raw observer is attached.
It assigns no observer value to a possibly nonadmitted auxiliary
port image at \(a_b\). Set \(R_b=kr_b>0\).
Positive \eqref{bp-tail:positive-inward-absolute}
and negative \eqref{bp-tail:negative-inward-absolute} monotonicity then
give a unique actual negative match for every \(a\ge a_b\):
\begin{equation}\label{bp-tail:whole-tail-matching}
 \eta_0^+(a)=\eta_0^-(r(a)),\qquad R_b\le R(a)<\mathfrak c.
\end{equation}
The order of the actual centre images therefore extends matching beyond
\(gku<d_0\). The absolute output estimate below holds on all
\(r\ge r_c\), without requiring the stronger bound \(R\ge R_b\).
For an input of the original \(I_\theta\), the through trajectory
already reaches both the centre and \(\Sigma_-\). Equality at that
same centre and scalar uniqueness identify the matched negative graph
with this trajectory. Its output-window and closing inverse-image
admission are inherited before an observer is evaluated. The auxiliary
formal shots do not enlarge those observer domains.

We now bound the actual negative output gain at the fixed raw cut
\(t=d_q\), followed by its original regular section bridges.
Since \(R=kr\), the negative spatial radius is \(gR\).
In \(R=kr,W=k^2\eta\), with raw height \(z=g^2W\), the exact
positive-time field, after a positive constant time change, is
\[
 R'=g\widehat P(R)+Q_-(R)W,\qquad
 W'=W\{RB_-(R)-Cg^2(1-R)W\},\qquad
 Q_-(R)=\ell(g(1-R)).
\]
This retains the independent \(k\)-dependence of
\(\widehat P=k^2+\pi kR+3R^2-R^3\).
For the actual raw negative graph let \(t=g-v>0\),
\(\mathcal W=z/t^2\), \(b=gR\), and
\[
 \begin{gathered}
 p=3+\pi/r+r^{-2}-R=\widehat P(R)/R^2,\quad
 m_-(t)=3g-t+\pi kg^2/t+k^2g^3/t^2,\\
 Q_-(t)=\ell(g-t),\quad
 N_-(t,\mathcal W)=1+2g-t-Ct(g-t)\mathcal W,\quad
 D_-=m_-+Q_-\mathcal W.
 \end{gathered}
\]
Its original lip is \(\mathcal W_l=\ell_0gp\).
Corollary~\ref{deg:cor-negative-absolute-outer} proves actual passage
and original-buffer admission on all \(r\ge r_c,0<R<\mathfrak c\),
including \(p\to0\), by its inverse graph and strip climb.

Set \(\nu=6g/(1+2g)\), \(\lambda_g=2\nu\). The exact density
\eqref{deg:quadratic-outer-density} on this negative graph is
\begin{equation}\label{bp-tail:negative-exact-density}
 \begin{gathered}
 I_- =\int_{\log z_l}^{\log z_T}J_-\,d\log z,\\
 J_-=-\frac{t}{N_-}
 \left[3-\nu+\{-A+(3-\nu)C(g-t)\}\mathcal W
                                      -\frac{\pi kg^2}{t^2}\right].
 \end{gathered}
\end{equation}
The output gain contains \(e^{I_-}\), so an upper bound for
\(I_-\) suffices. Choose \(0<\rho_0<\inf\mathfrak c/4\).
For \(R\le\rho_0\), \(p\) has bounds and reciprocal bounds on
\(r\ge r_c\). Its weak \(\log z\)-length is \(O(1+|\log g|)\),
and \(|J_-|\le Cg\) there. This part contributes \(O(g(1+|\log g|))\).
For \(R\ge\rho_0\), choose a fixed small auxiliary height
\(\omega>0\) with
\(\sup|{-A+(3-\nu)C(g-t)}|\omega\le1/4\), then a fixed
\[
             k_{\rm tail}\le\frac{\rho_0^2}{4\Pi}.
\]
It is a separate quadratic-chart bound, fixed before any final raw
coefficient shrink. While \(\mathcal W\le\omega\),
\(t\ge b=gR\) gives \(\pi kg^2/t^2\le\Pi k/\rho_0^2\le1/4\).
After choosing \(g_{\rm tail}\) with \(\nu<1/4\), the bracket
in \eqref{bp-tail:negative-exact-density} is at least \(2\).
Hence \(J_-\le0\) throughout the arbitrarily long root dwell.
No logarithm of \(p\) is estimated by a coefficient bound.
Choose a fixed inner cut \(d_*\ll\omega\) and then shrink \(g\);
on the strong part up to \(d_*\), \(\mathcal W\ge\omega\) and
\(D_-\ge c\mathcal W\). Thus \(d\log z\le C\,d\log t\).
The remaining density
has majorant \(C(t+kg^2/t)\), with integral at most
\(C(d_*+kg^2/b)\le C(d_*+g/r_c)\).
The compact strong continuation from this auxiliary \(d_*\) to the
original fixed raw cut \(d_q\) contributes a fixed bound. Thus on the whole original
negative interval
\begin{equation}\label{bp-tail:negative-density-value}
                              I_-\le C.
\end{equation}
The singular root dwell is controlled by its sign; this estimate imposes
no condition on \(g\log k\).

Here \(T\) denotes \(t=d_q\). Let \(\psi_-\) include the original
regular bridge from this cut to \(\Sigma_-\) and the divided closing,
and choose \(c_-=\psi_-'(z_{\rm ref})>0\) at a coefficient-only
buffered reference. It and its reciprocal are observer units.
The \(s=-1\) specialization of
\eqref{deg:quadratic-outer-exact-slope}, valid algebraically on this
whole negative domain, and the exact label change give
\begin{equation}\label{bp-tail:negative-full-gain}
 \Gamma_{-,b_R}'=c_-g^2\ell_0^{-1-\nu}g^{-\nu}(gR)^{-\lambda_g}
                    p^{-\nu}\mathcal K_-(r),\qquad
 b_R=-R^2/2,\quad0<\mathcal K_-\le C.
\end{equation}
Here \(\mathcal K_-\) is the product \(U_s\) of that formula,
without its initial \(p^{-\nu}\); its upper bound follows from
\eqref{bp-tail:negative-density-value} and the remaining unit factors.
Here \(\mathscr B_-\) is the original moving-initial transverse
factor of \eqref{deg:quadratic-outer-moving-factor}, a negative unit.
The terminal factors are evaluated on their admitted original images.
The power \(p^{-\nu}\), the exact \(\lambda_g\), and the absolute
gain are all retained; \(\mathcal K_-^{-1}\) is not needed.
Changing back to \(b_r=-r^2/2\) gives exactly
\(\Gamma_{-,b_r}'=k^2\Gamma_{-,b_R}'\); no \(k^{-\lambda_g}\)
has been declared a unit.
Using the actual center variation \eqref{bp-tail:negative-inward-absolute},
\begin{equation}\label{bp-tail:negative-consumed-gain}
 0<\frac{d\Gamma_-}{d\log\eta_0^-}
   =\frac{g\widehat P(R)}{B_-\Lambda_-}\Gamma_{-,b_R}'
   \le Cc_-g^3g^{-\nu-\lambda_g}R^{2-\lambda_g}p^{1-\nu}
   \le Cc_-g^3.
\end{equation}
Indeed \(0<p\le C\) on \(r\ge r_c\), \(R<\mathfrak c\), and
\(2-\lambda_g>1,1-\nu>0\). These positive powers pay both
\(R\to0\) and \(p\to0\). Also \(g^{-\nu-\lambda_g}\) is a
unit because \(\nu+\lambda_g=O(g)\), independently of \(k\).
Thus differentiation with respect to the actual centre height absorbs
the possibly singular root factor in the output gain.

Choose \(c_+>0\) as the analogous ordinary input observer slope unit.
Equations~\eqref{bp-tail:positive-full-gain} and
\eqref{bp-tail:positive-inward-absolute} give
\[
 \frac{d\Gamma_+}{d\log\eta_0^+}
   =\frac{-\Gamma_{+,a}'M(a)}{aB_0(a)\Lambda_+(a)}
                           \ge c(d_0)c_+a^3.
\]
On the same actual matched orbit the complete physical derivative is
therefore the exact ratio
\[
 0<\Pi'(x)=
 \frac{d\Gamma_-/d\log\eta_0^-}{d\Gamma_+/d\log\eta_0^+}
 \le C\frac{c_-}{c_+}\frac{g^3}{a^3}\le C(d_0)g^3.
\]
Only the ratio of the original regular observer units is bounded here;
the full exterior gains and their endpoint powers remained in
\eqref{bp-tail:positive-full-gain} and \eqref{bp-tail:negative-full-gain}.
The input center subtraction \(a_\theta\) has zero phase derivative.
The ordinary \(C^3\) divided-closing chain rule of
Lemma~\ref{deg:lem-scalar-moving-identities} is uniform through
\(m=0\) on these original buffered images. All coefficients remain
frozen, independently of \(g,k\), the ghost gap and the phase scale.
\end{proof}


\begin{corollary}[The bounded positive quadratic carrier]
\label{deg:cor-positive-quadratic-whole-carrier}
Fix \(\Pi_0<\infty\), the original nested raw-section windows,
and their ordinary \(C^3\) regular observers and divided finite
closing. The passive compact \((A,C)\)-set is contained in their
prescribed common coefficient neighbourhood, with the original
buffer and observer margins. Use fibres with \(K_{\rm fin}>0\).
For the field \eqref{deg:quadratic-field} there are independent
upper bounds for \(g,k>0\), one common multiplier bound
\(0<m\le m_*\le m_0\), and a fixed raw central-height cap
\(H_{\rm cap}>0\) with the following property.
On each retained real, fold or adjacent ghost fibre,
\(0<\pi\le\Pi_0\), the original enlarged windows admit one
connected physical carrier below \(H_{\rm cap}\). Its lower
interval has at most one curvature zero, which is simple and
crosses downward in the increasing physical input. The actual
\(F(x)=x-m\Pi_\theta(x)\) has at most two zeros of \(F'\), counted
with multiplicity, on the entire carrier. Its overlapping upper
complement has \(F'>0\). The ghost raw tail \(gku\ge d_0\) also
has \(F'>1/2\); it is retained on the original carrier together
with every intervening input. On the remaining root-free fibres,
\(F'>0\).
\end{corollary}

\begin{proof}
We combine the low-height curvature orientation with the paired upper
bound and the raw-tail contraction, all in the original affine input.
The strict overlaps below ensure that these estimates control every
point of the connected carrier before Rolle's theorem is applied.

Use the original affine input and divided output of
Lemma~\ref{lem:closing} and \eqref{sharp:upper-affine}.
Make the comparison and reserve choices of
Proposition~\ref{deg:prop-positive-quadratic}: \(K_0,r_c,X,\Delta_0,\pi_\flat\),
then \(\delta_q,\delta,M\), with \(M\delta\gg\Pi_0\).
Prepare the bounded supplier on \(u<4M\) and retain \(u\le3M\).
Choose \(d_0\) for the small error in
\eqref{deg:quadratic-large-S-curvature}, with
\(0<d_{\rm in}<d_0<\min(a_0,d_q/2)\).
Fix \(0<H_{\rm cap}<z_*\) before coefficient and multiplier bounds,
and apply Proposition~\ref{bp-tail:prop-actual-raw-tail} with
\(a_b=d_0/2\), \(h_c=H_{\rm cap}/2\); include
\(g_{\rm tail},k_{\rm tail}\) in the final rectangle.

The same-centre suppliers give \(N_\pm\ge1/2\), and
\[
 Z_-(r)\downarrow0\quad(kr\uparrow\mathfrak c),\qquad
 Z_-([r_c,\mathfrak c/k))=(0,Z_-(r_c)].
\]
The comparison action margin places \(Z_+(u_{\rm th})\) inside
this image and, by its bounded action, above \(Z_-(r_{\max})\)
for a fixed \(r_{\max}>r_c\). Strict monotonicity therefore gives
every retained lip its unique actual match,
\[
 Z_+(u)=Z_-(r),\qquad
 R_u=\frac{uB_+\widehat P(R)}{RB_-P(u)}\frac{N_+}{N_-}>0.
\]
Real/fold inputs have \(u_{\rm th}\le u<a\), with only \(u<d\)
at the fold; ghost curvature inputs use the original enlarged
raw-window preimage with \(u\ge u_{\rm th},gku<d_0\).
The separate outward suppliers admit both raw images before observers.
Scalar order and Lemma~\ref{lem:loc-through} make the window
preimages and original speed, height and stop clips intervals.

The comparison is admitted to the mother interval: its central
raw height is below \(h_c\), and its central speed is below the
original margin after shrinking coefficients. Below it the positive
port has \(Y/P>\ell_0\) and the centre has \(Y/P<\ell_0\).
Ghost positivity gives an intermediate lip hit. At every hit the
action-clock derivative is strictly negative, so the hit is unique.
Its radial derivative has the opposite sign, whereas scalar first
variation in the original input is positive. Implicit differentiation
therefore gives \(a_x<0\) in that same affine input. This first-hit
construction and the centre-image identification above preserve both
original output and closing clips. Allocate
\[
 I_{\rm tail}=\{x\in I_\theta:x\le x_{\rm comp},\ \Delta<0,
                                      \ a(x)\ge d_0\},\qquad
 I_{\rm low}=\{x\in I_\theta:x\le x_{\rm comp}\}\setminus I_{\rm tail}.
\]
The tail is an initial interval or empty. Its entry bound
\(H\le(gk)^2Cg e^{-c/g}<h_c\) puts its whole admitted part in
Proposition~\ref{bp-tail:prop-actual-raw-tail}, including \(a=d_0\).
That estimate strictly overlaps the native curvature interval on
\(a_b\le a<d_0\); the connected \(I_{\rm low}\) retains the
strict hypothesis \(a<d_0\).

On the bounded range, the two action-three supplier bounds and
bounded rows of \(H_\pm\) give, by the separated quotient rule,
\begin{equation}\label{deg:quadratic-bounded-product-rows}
 \sum_{i+j\le2}|\mathcal D_+^i\mathcal D_-^j(n-1)|\le C_Mg,
 \qquad \mathcal D_+\mathcal D_-n
       =-\frac{(\mathcal D_+N_+)(\mathcal D_-N_-)}{N_-^2}.
\end{equation}
Thus \eqref{deg:quadratic-bounded-lip-curvature} consumes the full
product and section rows. Proposition~\ref{deg:prop-positive-quadratic}
gives the downward crossings on the bounded lip, fixed-normal and
joint-growing packages. The exact resection
\eqref{deg:joint-root-exact-lip-resection} and composition
\eqref{deg:joint-root-physical-composition} identify the same map
on their overlaps, retaining coefficient memory and both gains.

For \(\sigma_1=\pi_\flat/(24M)\), use the positive bounds
\(R_0(\delta),R_0(\sigma_1)\) from
\eqref{deg:quadratic-large-S-action}, and choose
\[
 k_0<\frac{\min\{R_0(\delta),R_0(\sigma_1)\}}{4M},\qquad
 k_0\le\min\{\delta/(2M),k_{\rm tail}\}.
\]
These enforce the proposition's mixed-range exclusion; the remaining
\(R\ge\delta,S\le\delta\) and \(S\ge\delta\) packages have positive
curvature. The closed alternatives cover \(r\ge2M,u\ge2M\),
within the enlarged domains \(S,R<\delta_q\) and
\(S>\delta/2,a<2d_0\); the bounded range \(u<4M\) overlaps at
\(2M\le u\le3M\). Every equality face is retained.
Hence every curvature zero on the connected lower interval is a
simple downward crossing, so there is at most one. Rolle applied
to \(F''=-m\Pi_\theta''\) gives at most two zeros of \(F'\),
with multiplicity; a zero at the curvature crossing is double.

The paired lemma gives \(0<\Pi_\theta'\le C_*\) on the original
buffer interval \([\eta_{\rm comp},\eta_{\max}]\), with
\(\eta_{\rm comp}=gZ_+(u_*)<gZ_+(u_{\rm th})<\eta_{\max}\).
The signed high package applies for
\(\eta_0>gK_H(1+g\pi^2)\), \(H=(gk)^2\eta_0<z_*\).
Shrink independently so
\[
 gK_H(1+g\Pi_0^2)<\eta_{\max}/2,\qquad
 (gk)^2\eta_{\max}<H_{\rm cap}/2.
\]
The paired and high packages strictly overlap, retain the high equality
in the paired package and cover every intervening upper input under
the original clips. Finally take
\(m_*\max(C_*,C_H,C(d_0)g_0^3,C_{\rm contraction})<1/2\).
The tail and upper complement have \(F'>1/2\), adding no critical
point to the same physical \(F\). For \(\pi\le\pi_\flat\),
Lemma~\ref{deg:lem-raw-jacobian} gives \(F'>0\).
\end{proof}
